\documentclass[11pt]{article}
\usepackage{amsmath}
\usepackage{graphicx}
\usepackage[margin=1in]{geometry}
\usepackage{hyperref}
\usepackage{minted}
%\usepackage{tocloft}
%\usepackage{inconsolata}
%\usepackage{xcolor}
%\usepackage{etoolbox}
%\usepackage{fancyvrb}


%\definecolor{bg}{rgb}{0.95,0.95,0.95}

% Make the background color of minted listings this color
\setminted{bgcolor=bg}

% Add a frame around the listing
\setminted{frame=lines, framerule=1pt}

% Add a bit of space between the frame and the code
\setminted{framesep=5pt}

% Make the font size a bit smaller
\setminted{fontsize=\small}

\RecustomVerbatimEnvironment{Verbatim}{Verbatim}{fontsize=\small, frame=lines, framerule=1pt, framesep=5pt}


\title{VibeVolts: A Python-Based Toolkit for Space Environment Simulation}
\author{Jules, AI Software Engineer}
\date{\today}

\begin{document}

\maketitle

\begin{abstract}
VibeVolts is a Python-based simulation toolkit for space environment modeling. It provides a set of tools to initialize, propagate, and analyze the state of various space-based and ground-based assets. Key features include orbit propagation, exclusion analysis, and 3D visualization. The toolkit is designed with a modular, performance-oriented architecture that heavily leverages NumPy for vectorized operations.
\end{abstract}

\tableofcontents

\section{Introduction}
VibeVolts is a fast and efficient tool for simulating the visibility of space-based assets. It has been refactored into a modular structure with a focus on performance, leveraging established astronomical libraries like \texttt{astropy}. The core capabilities of the toolkit include satellite propagation from TLE data, celestial body positioning, viewing exclusion analysis (from the Sun, Moon, and Earth), and interactive 3D visualization of simulation scenarios using \texttt{plotly}.

\section{Core Data Structures}
The simulation state is managed within a single, large dictionary created by composition.

\subsection{The \texttt{simulation\_data} Dictionary}
This is the central, composable data structure that holds the entire simulation state. A minimal simulation is created with \texttt{create\_empty\_simulation} and then populated by other functions.

\begin{minted}{python}
{
    'start_time': datetime,
    'counts': {
        'celestial': 2,
        'satellites': num_satellites,
        'observatories': num_observatories,
        'red_satellites': num_red_satellites
    },
    'celestial': {
        'position': np.zeros((2, 3)),
        'velocity': np.zeros((2, 3)),
        'acceleration': np.zeros((2, 3)),
    },
    'satellites': {
        'position': np.zeros((num_satellites, 3)),
        'velocity': np.zeros((num_satellites, 3)),
        'acceleration': np.zeros((num_satellites, 3)),
        'orbital_elements': np.zeros((num_satellites, 6)),
        'epochs': [],
        'pointing': np.zeros((num_satellites, 3)),
        'pointing_state': np.zeros((num_satellites, 2), dtype=int),
        'detector': np.zeros((num_satellites, 7)),
    },
    'observatories': { ... },
    'red_satellites': { ... },
    'fixedpoints': {
        'position': np.zeros((num_points, 3)),
        'visibility': np.zeros((num_points, num_satellites), dtype=int)
    },
    'pointing_spheres': {},
    'delta_time': 60.0
}
\end{minted}

Key components include:
\begin{itemize}
    \item \texttt{counts}: Stores the number of entities in the simulation.
    \item \texttt{celestial}: Holds data for the Sun and Moon.
    \item \texttt{satellites}: Contains arrays for satellite properties like position, velocity, orbital elements, and detector characteristics.
    \item \texttt{fixedpoints}: Stores the positions of static observation targets and their visibility table.
    \item \texttt{pointing\_spheres}: Caches pre-computed pointing vector grids to avoid redundant calculations.
\end{itemize}

\subsection{Radiometric and Physical Data}
The \texttt{radiometry\_data.py} module provides the \texttt{FILTER\_DATA} dictionary, which contains standard data for various astronomical filters (e.g., Johnson-Cousins, SDSS, JWST). It also defines physical constants such as \texttt{AU\_M} (Astronomical Unit) and \texttt{RSUN\_M} (Solar Radius).

\section{Core Modules and Algorithms}

\subsection{\texttt{simulation.py} - Simulation Setup}
This module is responsible for creating the basic simulation structure.
\begin{itemize}
    \item \texttt{create\_empty\_simulation}: Initializes a minimal, empty data structure.
    \item \texttt{add\_celestial\_bodies}: Adds structures for the Sun and Moon.
    \item \texttt{add\_fixed\_points}: Generates and adds a set of static reference points.
\end{itemize}

\subsection{\texttt{propagation.py} - Orbit Dynamics}
This module handles the dynamics of the simulation.
\begin{itemize}
    \item \texttt{propagate\_satellites\_new}: Updates satellite positions using a vectorized Keplerian propagator.
    \item \texttt{celestial\_update}: Uses `astropy` to get high-precision positions for the Sun and Moon.
\end{itemize}

\subsection{\texttt{visibility.py} - Line-of-Sight Analysis}
This module determines what a satellite can "see."
\begin{itemize}
    \item \texttt{exclusion}: Checks if a satellite's line of sight is blocked by the Sun, Moon, or Earth.
    \item \texttt{solarexclusion}: A fully vectorized function to check for solar exclusion across all satellites.
    \item \texttt{update\_visibility\_table}: Iterates through all satellites and fixed points to build a visibility matrix.
\end{itemize}

\subsection{\texttt{pointing.py} - Satellite Pointing Control}
This module manages where the satellites are pointing.
\begin{itemize}
    \item \texttt{generate\_pointing\_sphere}: Creates and caches a pointing grid using a Fibonacci lattice.
    \item \texttt{update\_satellite\_pointing}: Updates each satellite's pointing vector based on its assigned grid and current index.
\end{itemize}

\section{Demonstrations and Usage}
The toolkit includes a suite of demo scripts (\texttt{demo*.py}) to showcase its features. Key demos include:
\begin{itemize}
    \item \texttt{demo1.py}: Basic propagation of a standard set of satellites.
    \item \texttt{demo\_exclusion\_table.py}: Calculates and displays the visibility matrix as a heatmap.
    \item \texttt{demo\_sky\_scan.py}: Simulates a sky scan from a GEO satellite to map celestial exclusion zones.
\end{itemize}
All demos can be run from a single script:
\begin{minted}{bash}
python all_demos.py
\end{minted}

\section{Dependencies}
The following Python libraries are required:
\begin{itemize}
    \item \texttt{numpy}
    \item \texttt{astropy}
    \item \texttt{jplephem}
    \item \texttt{sgp4}
    \item \texttt{plotly}
    \item \texttt{scipy}
    \item \texttt{ipython}
\end{itemize}
They can be installed with pip:
\begin{minted}{bash}
pip install numpy astropy jplephem sgp4 plotly scipy ipython
\end{minted}

\section{Complete Listings}
\subsection{\texttt{\_\_init\_\_.py}}
\begin{Verbatim}
# This file is intentionally left blank.
\end{Verbatim}
\subsection{\texttt{all\_demos.py}}
\begin{Verbatim}
import plotly.graph_objects as go
from IPython.display import display, HTML


# Import all the demo functions
from demo1 import demo1
from demo2 import demo2
from demo3 import demo3
from demo4 import demo4
from demo_fixedpoints import demo_fixedpoints
from demo_exclusion_table import demo_exclusion_table
from demo_pointing_plot import demo_pointing_plot

from demo_lambertian import demo_lambertian
from demo_sky_scan import demo_sky_scan
from demo_pointing_vectors import demo_pointing_vectors
from demo_pointing_sequence import demo_pointing_sequence
from demo_constellation import demo_constellation
from show_geo_search import show_geo_search


def run_all_demos(save_html=False):
    """
    Runs all demo functions, and either shows them inline or saves them to a single HTML file.

    Args:
        save_html (bool): If True, saves plots to HTML. If False, displays plots inline.
    """
    demo_functions = [
 #       demo1,
  #      demo2,
  #      demo3,
  #      demo4,
  #      demo_fixedpoints,
  #      demo_exclusion_table,
  #      demo_pointing_plot,
  #      demo_lambertian,
  #      demo_sky_scan,
  #      demo_pointing_vectors,
  #      demo_pointing_sequence,
        demo_constellation,
        show_geo_search,
    ]

    figs = [] # This is a list
    print("--- Running All Demos ---")
    for func in demo_functions:
        print(f"\n\033[91m--- Executing {func.__name__} ---\033[0m")
        res = func()
        if isinstance(res, go.Figure):
            figs.append(res)
            print("appending a single item")
        elif isinstance(res, tuple):
            print("I have a tuple")
            for fig in res:
                figs.append(fig)
                print("extracting an item from a tuple")

    if save_html:
        with open("all_demo_plots.html", "w") as f:
            f.write("<html><head><title>VibeVolts Demos</title></head><body>")
            f.write("<h1>VibeVolts Demo Plots</h1>")
            for i, fig in enumerate(figs):
                f.write(f"<h2>Demo {i+1}: {fig.layout.title.text if fig.layout.title else ''}</h2>")
                f.write(fig.to_html(full_html=False, include_plotlyjs='cdn'))
            f.write("</body></html>")
        print("\n--- All demos complete. Plots saved to all_demo_plots.html ---")
    else:
        print("\n--- Displaying all demo plots inline ---")
        print(len(figs))
        for fig in figs:
           fig.show()

if __name__ == '__main__':
    run_all_demos()
\end{Verbatim}
\subsection{\texttt{constants.py}}
\begin{Verbatim}
# --- Global Constants for Array Indices ---
# These constants define column indices for numpy arrays, making the
# code more readable and preventing errors from using "magic numbers".

# -- Radii in Meters --
EARTH_RADIUS = 6378137.0
MOON_RADIUS = 1737400.0

# -- Detector Array Indices --
DETECTOR_APERTURE_IDX = 0      # Aperture size in meters
DETECTOR_PIXEL_SIZE_IDX = 1    # Pixel size in radians
DETECTOR_QE_IDX = 2            # Quantum efficiency as a fraction (0.0 to 1.0)
DETECTOR_PIXELS_IDX = 3        # Total number of pixels in the detector (count)
DETECTOR_SOLAR_EXCL_IDX = 4    # Solar exclusion angle in radians
DETECTOR_LUNAR_EXCL_IDX = 5    # Lunar exclusion angle in radians
DETECTOR_EARTH_EXCL_IDX = 6    # Earth exclusion angle (above the limb) in radians

# -- Orbital Elements Array Indices --
ORBITAL_A_IDX = 0              # Semi-major axis in meters
ORBITAL_E_IDX = 1              # Eccentricity (dimensionless)
ORBITAL_I_IDX = 2              # Inclination in radians
ORBITAL_RAAN_IDX = 3           # Right Ascension of the Ascending Node in radians
ORBITAL_ARGP_IDX = 4           # Argument of Perigee in radians
ORBITAL_M_IDX = 5              # Mean Anomaly in radians

# -- Pointing State Array Indices --
POINTING_COUNT_IDX = 0         # Number of points in the pointing grid
POINTING_PLACE_IDX = 1         # Current index in the pointing sequence
\end{Verbatim}
\subsection{\texttt{constellation.py}}
\begin{Verbatim}
import numpy as np
from datetime import datetime
from typing import Dict, Any
import math
import random

from constants import (
    ORBITAL_A_IDX, ORBITAL_E_IDX, ORBITAL_I_IDX,
    ORBITAL_RAAN_IDX, ORBITAL_ARGP_IDX, ORBITAL_M_IDX,
    POINTING_COUNT_IDX, POINTING_PLACE_IDX
)
from pointing import generate_pointing_sphere

def geos(sim_data, n,  fov) -> None:
    """
    Creates n equally spaced satellites in GEO and adds them to the 'satellites' group in the simulation.

    Args:
        sim_data: The main simulation data dictionary.
        n: The number of satellites to create.
        fov: The diameter of the field of view of the satellite in radians.
    """
    # Calculate solid angle
    theta = fov / 2
    solid_angle = 2 * np.pi * (1 - np.cos(theta))

    # Calculate grid_points - blow things up by 0.25 for overlap
    grid_points = int(4 * np.pi / solid_angle * 1.25)

    # Generate and store the pointing sphere and place in ['pointing_sphers'][n]
    generate_pointing_sphere(sim_data, grid_points)

    orbital_elements_list = []
    epochs_list = []
    pointing_state_list = []

    # Geostationary orbit semi-major axis in meters
    a = 42164000.0

    # Create a set of elements evenly spaced around the equator in
    # orbital_elemens_list
    for i in range(n):
        elements = np.zeros(6)
        elements[ORBITAL_A_IDX] = a
        elements[ORBITAL_E_IDX] = 0.0
        elements[ORBITAL_I_IDX] = 0.0
        elements[ORBITAL_RAAN_IDX] = i * 2 * np.pi / n
        elements[ORBITAL_ARGP_IDX] = 0.0
        elements[ORBITAL_M_IDX] = 0.0
        orbital_elements_list.append(elements)
        epochs_list.append(sim_data['start_time'])

   # And these are the pointing state of the satellite. Make the position random.

        pointing_state = np.zeros(2, dtype=int)
        pointing_state[POINTING_COUNT_IDX] = grid_points
        pointing_state[POINTING_PLACE_IDX] = random.randint(0,grid_points-1)

        pointing_state_list.append(pointing_state)

    orbital_elements = np.array(orbital_elements_list, dtype=float)
    pointing_state_array = np.array(pointing_state_list, dtype=int)

    if 'satellites' not in sim_data:
        sim_data['counts']['satellites'] = n
        sim_data['satellites'] = {
            'position': np.zeros((n, 3), dtype=float),
            'velocity': np.zeros((n, 3), dtype=float),
            'acceleration': np.zeros((n, 3), dtype=float),
            'orbital_elements': orbital_elements,
            'epochs': epochs_list,
            'pointing': np.zeros((n, 3), dtype=float),
            'pointing_state': pointing_state_array,
            'detector': np.zeros((n, 7), dtype=float),
        }
    else:
        # Append to existing satellites
        sim_data['counts']['satellites'] += n
        sim_data['satellites']['position'] = np.vstack([sim_data['satellites']['position'],
np.zeros((n, 3), dtype=float)])
        sim_data['satellites']['velocity'] = np.vstack([sim_data['satellites']['velocity'],
np.zeros((n, 3), dtype=float)])
        sim_data['satellites']['acceleration'] = np.vstack([sim_data['satellites']['acceleration'],
np.zeros((n, 3), dtype=float)])
        sim_data['satellites']['orbital_elements'] = np.vstack([sim_data['satellites']['orbital_elements'],
orbital_elements])
        sim_data['satellites']['epochs'].extend(epochs_list)
        sim_data['satellites']['pointing'] = np.vstack([sim_data['satellites']['pointing'],
np.zeros((n, 3), dtype=float)])
        sim_data['satellites']['pointing_state'] = np.vstack([sim_data['satellites']['pointing_state'],
pointing_state_array])
        sim_data['satellites']['detector'] = np.vstack([sim_data['satellites']['detector'],
np.zeros((n, 7), dtype=float)])
\end{Verbatim}
\subsection{\texttt{demo1.py}}
\begin{Verbatim}
from datetime import datetime, timezone, timedelta
import plotly.graph_objects as go

from demo_common import initialize_standard_simulation
from propagation import celestial_update, propagate_satellites_new
from plotting_3d import plot_3d_scatter

def demo1() -> go.Figure:
    """
    Runs a full demonstration of the simulation tools.

    This function generates and returns a Plotly figure object but does not
    display it. The caller is responsible for rendering the plot.

    Returns:
        The Plotly figure object for the satellite positions plot.
    """
    sim_start_time = datetime(2025, 7, 27, 22, 27, 0, tzinfo=timezone.utc)
    sim_data = initialize_standard_simulation(sim_start_time)

    # --- Celestial Body Updates ---
    sim_data = celestial_update(sim_data, sim_start_time)
    print("\n--- Celestial Positions at Start Time ---")
    print(f"Time: {sim_start_time.isoformat()}")
    print(sim_data['celestial']['position'])

    # --- Satellite Propagation and Plotting ---
    time_t1 = sim_start_time + timedelta(hours=1, minutes=30)
    print(f"\n--- Propagating satellites to T1: {time_t1.isoformat()} ---")
    sim_data = propagate_satellites_new(sim_data, time_t1)

    fig = plot_3d_scatter(
        positions=sim_data['satellites']['position'],
        title=f"Satellite Positions at {time_t1.isoformat()}",
        plot_time=time_t1,
        marker_size=2,
        trace_name='Satellites'
    )
    return fig
\end{Verbatim}
\subsection{\texttt{demo2.py}}
\begin{Verbatim}
import numpy as np
from datetime import datetime, timezone, timedelta
import plotly.graph_objects as go

from demo_common import initialize_standard_simulation
from propagation import celestial_update, propagate_satellites_new

def demo2() -> go.Figure:
    """
    Runs a demonstration plotting satellite and celestial positions.

    This function generates and returns a Plotly figure object but does not
    display it. The caller is responsible for rendering the plot.

    Returns:
        The Plotly figure object.
    """
    sim_start_time = datetime(2025, 7, 27, 22, 27, 0, tzinfo=timezone.utc)
    sim_data = initialize_standard_simulation(sim_start_time)

    # --- Satellite and Celestial Propagation ---
    positions_t0 = sim_data['satellites']['position'].copy()
    sim_data = celestial_update(sim_data, sim_start_time)
    celestial_pos_t0 = sim_data['celestial']['position'].copy()

    time_t1 = sim_start_time + timedelta(seconds=300)
    sim_data = propagate_satellites_new(sim_data, time_t1)
    positions_t1 = sim_data['satellites']['position'].copy()
    sim_data = celestial_update(sim_data, time_t1)
    celestial_pos_t1 = sim_data['celestial']['position'].copy()

    # --- 3D Plotting of All Time Steps ---
    print("\n--- Generating 3D plot for Demo 2 ---")
    earth_radius = 6378137.0
    vector_scale = 10 * earth_radius
    fig = go.Figure()

    fig.add_trace(go.Scatter3d(
        x=positions_t0[:, 0], y=positions_t0[:, 1], z=positions_t0[:, 2],
        mode='markers', marker=dict(size=5, color='blue'), name='Sats (T=0s)'
    ))
    fig.add_trace(go.Scatter3d(
        x=positions_t1[:, 0], y=positions_t1[:, 1], z=positions_t1[:, 2],
        mode='markers', marker=dict(size=5, color='red'), name='Sats (T=300s)'
    ))

    sun_vec_t0 = celestial_pos_t0[0] / np.linalg.norm(celestial_pos_t0[0]) * vector_scale
    moon_vec_t0 = celestial_pos_t0[1] / np.linalg.norm(celestial_pos_t0[1]) * vector_scale
    fig.add_trace(go.Scatter3d(x=[0, sun_vec_t0[0]], y=[0, sun_vec_t0[1]], z=[0, sun_vec_t0[2]],
mode='lines', line=dict(color='orange', width=4), name='Sun (T=0s)'))
    fig.add_trace(go.Scatter3d(x=[0, moon_vec_t0[0]], y=[0, moon_vec_t0[1]], z=[0, moon_vec_t0[2]],
mode='lines', line=dict(color='gray', width=4), name='Moon (T=0s)'))

    sun_vec_t1 = celestial_pos_t1[0] / np.linalg.norm(celestial_pos_t1[0]) * vector_scale
    moon_vec_t1 = celestial_pos_t1[1] / np.linalg.norm(celestial_pos_t1[1]) * vector_scale
    fig.add_trace(go.Scatter3d(x=[0, sun_vec_t1[0]], y=[0, sun_vec_t1[1]], z=[0, sun_vec_t1[2]],
mode='lines', line=dict(color='yellow', width=4), name='Sun (T=300s)'))
    fig.add_trace(go.Scatter3d(x=[0, moon_vec_t1[0]], y=[0, moon_vec_t1[1]], z=[0, moon_vec_t1[2]],
mode='lines', line=dict(color='lightgray', width=4), name='Moon (T=300s)'))

    u_sphere = np.linspace(0, 2 * np.pi, 100)
    v_sphere = np.linspace(0, np.pi, 100)
    x_earth = earth_radius * np.outer(np.cos(u_sphere), np.sin(v_sphere))
    y_earth = earth_radius * np.outer(np.sin(u_sphere), np.sin(v_sphere))
    z_earth = earth_radius * np.outer(np.ones(np.size(u_sphere)), np.cos(v_sphere))
    fig.add_trace(go.Surface(x=x_earth, y=y_earth, z=z_earth, colorscale='Blues',
showscale=False, opacity=0.5, name='Earth'))

    fig.update_layout(
        title=f"Satellite and Celestial Positions at 0 and 300 seconds",
        scene=dict(xaxis_title='X (m)', yaxis_title='Y (m)', zaxis_title='Z (m)', aspectmode='data'),
        margin=dict(r=20, b=10, l=10, t=40)
    )

    return fig
\end{Verbatim}
\subsection{\texttt{demo3.py}}
\begin{Verbatim}
import numpy as np
from datetime import datetime, timezone, timedelta
import plotly.graph_objects as go

from propagation import add_satellites_from_tle, propagate_satellites_new
from simulation import create_empty_simulation

def demo3() -> go.Figure:
    """
    Runs a demonstration plotting a single LEO satellite trajectory.

    This function generates and returns a Plotly figure object but does not
    display it. The caller is responsible for rendering the plot.

    Returns:
        The Plotly figure object.
    """
    sim_start_time = datetime(2025, 7, 27, 22, 27, 0, tzinfo=timezone.utc)

    tle_data = '''LEO-TRAJECTORY
1 90201U 25005A   25210.50000000  .00000000  00000-0  00000-0 0  9991
2 90201  51.6000  10.0000 0010000  90.0000  20.0000 15.50000000    11
'''
    dummy_tle_path = "dummy_tle3.txt"
    with open(dummy_tle_path, "w") as f:
        f.write(tle_data)

    print(f"\n--- Starting Demo 3 ---")

    sim_data = create_empty_simulation(sim_start_time)
    add_satellites_from_tle(sim_data, dummy_tle_path, 'satellites')

    print(f"Initializing structures for {sim_data['counts']['satellites']} LEO satellite.")

    positions_over_time = []
    time_steps = np.arange(0, 91, 10)

    for minutes in time_steps:
        prop_time = sim_start_time + timedelta(minutes=int(minutes))
        sim_data = propagate_satellites_new(sim_data, prop_time)
        positions_over_time.append(sim_data['satellites']['position'][0])

    positions_array = np.array(positions_over_time)

    print("\n--- Generating 3D plot for Demo 3 ---")
    earth_radius = 6378137.0
    fig = go.Figure()

    fig.add_trace(go.Scatter3d(
        x=positions_array[:, 0], y=positions_array[:, 1], z=positions_array[:, 2],
        mode='lines', line=dict(color='red', width=4), name='Trajectory'
    ))

    fig.add_trace(go.Scatter3d(
        x=positions_array[:, 0], y=positions_array[:, 1], z=positions_array[:, 2],
        mode='markers', marker=dict(size=5, color='blue'),
        text=[f'T={t} min' for t in time_steps], hoverinfo='text', name='Time Steps'
    ))

    u_sphere = np.linspace(0, 2 * np.pi, 100)
    v_sphere = np.linspace(0, np.pi, 100)
    x_earth = earth_radius * np.outer(np.cos(u_sphere), np.sin(v_sphere))
    y_earth = earth_radius * np.outer(np.sin(u_sphere), np.sin(v_sphere))
    z_earth = earth_radius * np.outer(np.ones(np.size(u_sphere)), np.cos(v_sphere))
    fig.add_trace(go.Surface(x=x_earth, y=y_earth, z=z_earth, colorscale='Blues',
showscale=False, opacity=0.5, name='Earth'))

    fig.update_layout(
        title=f"Single LEO Satellite Trajectory over 90 Minutes",
        scene=dict(xaxis_title='X (m)', yaxis_title='Y (m)', zaxis_title='Z (m)', aspectmode='data'),
        margin=dict(r=20, b=10, l=10, t=40),
        legend_title_text='Trace'
    )

    return fig
\end{Verbatim}
\subsection{\texttt{demo4.py}}
\begin{Verbatim}
import numpy as np
from datetime import datetime, timezone, timedelta
import plotly.graph_objects as go

from propagation import add_satellites_from_tle, propagate_satellites_new
from simulation import create_empty_simulation

def demo4() -> go.Figure:
    """
    Runs a demonstration plotting a single GEO satellite trajectory.

    This function generates and returns a Plotly figure object but does not
    display it. The caller is responsible for rendering the plot.

    Returns:
        The Plotly figure object.
    """
    sim_start_time = datetime(2025, 7, 27, 22, 27, 0, tzinfo=timezone.utc)

    tle_data = '''GEO-TRAJECTORY
1 90301U 25006A   25210.50000000  .00000000  00000-0  00000-0 0  9991
2 90301   0.0500  45.0000 0001000  90.0000  20.0000  1.00270000    11
'''
    dummy_tle_path = "dummy_tle4.txt"
    with open(dummy_tle_path, "w") as f:
        f.write(tle_data)

    print(f"\n--- Starting Demo 4 ---")

    sim_data = create_empty_simulation(sim_start_time)
    add_satellites_from_tle(sim_data, dummy_tle_path, 'satellites')

    print(f"Initializing structures for {sim_data['counts']['satellites']} GEO satellite.")

    positions_over_time = []
    time_steps = np.arange(0, 24, 1)

    for hours in time_steps:
        prop_time = sim_start_time + timedelta(hours=int(hours))
        sim_data = propagate_satellites_new(sim_data, prop_time)
        positions_over_time.append(sim_data['satellites']['position'][0])

    positions_array = np.array(positions_over_time)

    print("\n--- Generating 3D plot for Demo 4 ---")
    earth_radius = 6378137.0
    fig = go.Figure()

    fig.add_trace(go.Scatter3d(
        x=positions_array[:, 0], y=positions_array[:, 1], z=positions_array[:, 2],
        mode='lines', line=dict(color='purple', width=4), name='Trajectory'
    ))

    fig.add_trace(go.Scatter3d(
        x=positions_array[:, 0], y=positions_array[:, 1], z=positions_array[:, 2],
        mode='markers', marker=dict(size=5, color='orange'),
        text=[f'T={t} hr' for t in time_steps], hoverinfo='text', name='Time Steps'
    ))

    u_sphere = np.linspace(0, 2 * np.pi, 100)
    v_sphere = np.linspace(0, np.pi, 100)
    x_earth = earth_radius * np.outer(np.cos(u_sphere), np.sin(v_sphere))
    y_earth = earth_radius * np.outer(np.sin(u_sphere), np.sin(v_sphere))
    z_earth = earth_radius * np.outer(np.ones(np.size(u_sphere)), np.cos(v_sphere))
    fig.add_trace(go.Surface(x=x_earth, y=y_earth, z=z_earth, colorscale='Blues',
showscale=False, opacity=0.5, name='Earth'))

    fig.update_layout(
        title=f"Single GEO Satellite Trajectory over 23 Hours",
        scene=dict(xaxis_title='X (m)', yaxis_title='Y (m)', zaxis_title='Z (m)', aspectmode='data'),
        margin=dict(r=20, b=10, l=10, t=40),
        legend_title_text='Trace'
    )

    return fig
\end{Verbatim}
\subsection{\texttt{demo\_common.py}}
\begin{Verbatim}
import numpy as np
from datetime import datetime
from typing import Dict, Any
from astropy.coordinates import solar_system_ephemeris

# Adjusting import paths for the new project structure
from propagation import add_satellites_from_tle, propagate_satellites_new
from simulation import create_empty_simulation, add_celestial_bodies, add_fixed_points
from observatories import add_observatories

def initialize_standard_simulation(start_time: datetime) -> Dict[str, Any]:
    """
    Initializes a standard simulation with a predefined set of satellites.

    This function consolidates TLE data, initializes the main data structure,
    and propagates satellites to their initial positions at the specified start
    time. This ensures the returned simulation state is fully populated with
    positions and radially-outward pointing vectors.

    Args:
        start_time: The timezone-aware datetime object for the simulation start.

    Returns:
        The fully initialized and propagated simulation data dictionary.
    """
    # Consolidated TLE data from all demos
    tle_data = """ISS (ZARYA)
1 25544U 98067A   25209.52203988  .00012111  00000+0  22159-3 0  9991
2 25544  51.6412 254.9961 0006733  98.4322 261.6813 15.49493393462383
NOAA 19
1 33591U 09005A   25209.38959223  .00000100  00000+0  97987-4 0  9993
2 33591  99.1533 244.3362 0013327 101.3725 258.7562 14.12510122810029
HST
1 20580U 90037B   25208.83160218  .00000113  00000+0  35999-4 0  9990
2 20580  28.4695 177.8391 0001259 138.5273 221.5822 15.09326468 23453
GEO-01
1 90001U 25001A   25209.50000000  .00000000  00000-0  00000-0 0  9991
2 90001   0.0500   0.0000 0001000   0.0000   0.0000  1.00270000    12
GEO-02
1 90002U 25001B   25209.50000000  .00000000  00000-0  00000-0 0  9992
2 90002   0.0500  36.0000 0001000   0.0000   0.0000  1.00270000    13
GEO-03
1 90003U 25001C   25209.50000000  .00000000  00000-0  00000-0 0  9993
2 90003   0.0500  72.0000 0001000   0.0000  45.0000  1.00270000    14
GEO-04
1 90004U 25001D   25209.50000000  .00000000  00000-0  00000-0 0  9994
2 90004   0.0500 108.0000 0001000   0.0000  90.0000  1.00270000    15
GEO-05
1 90005U 25001E   25209.50000000  .00000000  00000-0  00000-0 0  9995
2 90005   0.0500 144.0000 0001000   0.0000 135.0000  1.00270000    16
GEO-06
1 90006U 25001F   25209.50000000  .00000000  00000-0  00000-0 0  9996
2 90006   0.0500 180.0000 0001000   0.0000 180.0000  1.00270000    17
GEO-07
1 90007U 25001G   25209.50000000  .00000000  00000-0  00000-0 0  9997
2 90007   0.0500 216.0000 0001000   0.0000 225.0000  1.00270000    18
GEO-08
1 90008U 25001H   25209.50000000  .00000000  00000-0  00000-0 0  9998
2 90008   0.0500 252.0000 0001000   0.0000 270.0000  1.00270000    19
GEO-09
1 90009U 25001I   25209.50000000  .00000000  00000-0  00000-0 0  9999
2 90009   0.0500 288.0000 0001000   0.0000 315.0000  1.00270000    10
GEO-10
1 90010U 25001J   25209.50000000  .00000000  00000-0  00000-0 0  9990
2 90010   0.0500 324.0000 0001000   0.0000 360.0000  1.00270000    11
HEO-01 (MOLNIYA)
1 90011U 25002A   25209.50000000  .00000000  00000-0  00000-0 0  9995
2 90011  63.4000  50.0000 7500000  270.0000  45.0000  2.00560000    13
HEO-02 (MOLNIYA)
1 90012U 25002B   25209.50000000  .00000000  00000-0  00000-0 0  9996
2 90012  63.4000 110.0000 7500000  270.0000  90.0000  2.00560000    14
HEO-03 (MOLNIYA)
1 90013U 25002C   25209.50000000  .00000000  00000-0  00000-0 0  9997
2 90013  63.4000 170.0000 7500000  270.0000 135.0000  2.00560000    15
HEO-04 (MOLNIYA)
1 90014U 25002D   25209.50000000  .00000000  00000-0  00000-0 0  9998
2 90014  63.4000 230.0000 7500000  270.0000 180.0000  2.00560000    16
HEO-05 (MOLNIYA)
1 90015U 25002E   25209.50000000  .00000000  00000-0  00000-0 0  9999
2 90015  63.4000 290.0000 7500000  270.0000 225.0000  2.00560000    17
LEO-01 (POLAR)
1 90016U 25003A   25209.50000000  .00000000  00000-0  00000-0 0  9990
2 90016  98.0000 120.0000 0010000  90.0000  20.0000 14.50000000    18
LEO-02 (POLAR)
1 90017U 25003B   25209.50000000  .00000000  00000-0  00000-0 0  9991
2 90017  98.0000 240.0000 0010000  90.0000  40.0000 14.50000000    19
LEO-03
1 90018U 25003C   25209.50000000  .00000000  00000-0  00000-0 0  9992
2 90018  45.0000  80.0000 0010000  45.0000  60.0000 15.00000000    10
LEO-04
1 90019U 25003D   25209.50000000  .00000000  00000-0  00000-0 0  9993
2 90019  45.0000 200.0000 0010000  45.0000  80.0000 15.00000000    11
LEO-05
1 90020U 25003E   25209.50000000  .00000000  00000-0  00000-0 0  9994
2 90020  28.5000  10.0000 0010000  10.0000 100.0000 15.50000000    12
"""
    dummy_tle_path = "standard_tle.txt"
    with open(dummy_tle_path, "w") as f:
        f.write(tle_data)

    sim_data = create_empty_simulation(start_time)

    add_satellites_from_tle(sim_data, dummy_tle_path, 'satellites')
    add_observatories(sim_data, 0)
    add_celestial_bodies(sim_data)
    add_fixed_points(sim_data, 100)

    print(f"Initializing standard simulation with {sim_data['counts']['satellites']} satellites.")

    solar_system_ephemeris.set('jpl')

    print(f"Propagating satellites to start time: {start_time.isoformat()} to set initial state.")
    sim_data = propagate_satellites_new(sim_data, start_time)

    return sim_data
\end{Verbatim}
\subsection{\texttt{demo\_constellation.py}}
\begin{Verbatim}
from datetime import datetime, timezone, timedelta
import plotly.graph_objects as go

from demo_common import initialize_standard_simulation
from propagation import propagate_satellites_new
from plotting_3d import plot_3d_scatter
from constellation import geos

def demo_constellation() -> go.Figure:
    """
    Runs a demonstration of the constellation creation tools.

    This function generates and returns a Plotly figure object but does not
    display it. The caller is responsible for rendering the plot.

    Returns:
        The Plotly figure object for the satellite positions plot.
    """
    sim_start_time = datetime(2025, 7, 27, 22, 27, 0, tzinfo=timezone.utc)
    sim_data = initialize_standard_simulation(sim_start_time)

    num_sats_before = sim_data['counts']['satellites']
    # Create a constellation of 10 GEO satellites
    geos(sim_data, 10, 0.1)

    # Propagate the satellites for 1 hour
    time_t1 = sim_start_time + timedelta(hours=1)
    sim_data = propagate_satellites_new(sim_data, time_t1)

    fig = plot_3d_scatter(
        positions=sim_data['satellites']['position'][num_sats_before:],
        title=f"GEO Constellation at {time_t1.isoformat()}",
        plot_time=time_t1,
        marker_size=5,
        trace_name='GEO Satellites'
    )
    return fig

if __name__ == '__main__':
    fig = demo_constellation()
    fig.show()
\end{Verbatim}
\subsection{\texttt{demo\_exclusion\_debug\_print.py}}
\begin{Verbatim}
import numpy as np
from datetime import datetime, timezone

from demo_common import initialize_standard_simulation
from propagation import celestial_update
from visibility import update_visibility_table
from simulation import DETECTOR_SOLAR_EXCL_IDX, DETECTOR_LUNAR_EXCL_IDX, DETECTOR_EARTH_EXCL_IDX

def demo_exclusion_debug_print():
    """
    Demonstrates the debug printing feature of the exclusion function.

    This function runs the exclusion check for the first satellite against
    the first 100 fixed points and prints the detailed debug output for
    each check, as enabled by the `print_debug_for_sat` parameter.
    """
    print("\n--- Starting Demo: Exclusion Debug Print ---")
    sim_start_time = datetime(2025, 8, 1, 12, 0, 0, tzinfo=timezone.utc)
    sim_data = initialize_standard_simulation(sim_start_time)

    sim_data['satellites']['detector'][:, DETECTOR_SOLAR_EXCL_IDX] = np.deg2rad(30)
    sim_data['satellites']['detector'][:, DETECTOR_LUNAR_EXCL_IDX] = np.deg2rad(15)
    sim_data['satellites']['detector'][:, DETECTOR_EARTH_EXCL_IDX] = np.deg2rad(10)

    sim_data = celestial_update(sim_data, sim_start_time)

    print("\n--- Generating exclusion table for Satellite 0 vs First 100 Fixed Points (with debug print) ---")
    original_fixed_points = sim_data['fixedpoints']['position']
    sim_data['fixedpoints']['position'] = original_fixed_points[:100]

    update_visibility_table(sim_data, print_debug_for_sat=0)

    sim_data['fixedpoints']['position'] = original_fixed_points
    print("\n--- Debug Print Demo Complete ---")
\end{Verbatim}
\subsection{\texttt{demo\_exclusion\_table.py}}
\begin{Verbatim}
import numpy as np
from datetime import datetime, timezone
import plotly.graph_objects as go

from demo_common import initialize_standard_simulation
from propagation import celestial_update
from visibility import update_visibility_table
from simulation import DETECTOR_SOLAR_EXCL_IDX, DETECTOR_LUNAR_EXCL_IDX, DETECTOR_EARTH_EXCL_IDX

def demo_exclusion_table() -> go.Figure:
    """
    Demonstrates the creation and visualization of the exclusion table.

    This function generates and returns a Plotly figure object but does not
    display it. The caller is responsible for rendering the plot.

    Returns:
        The Plotly figure object.
    """
    print("\n--- Starting Demo: Exclusion Table Visualization ---")
    sim_start_time = datetime(2025, 8, 1, 12, 0, 0, tzinfo=timezone.utc)
    sim_data = initialize_standard_simulation(sim_start_time)

    sim_data['satellites']['detector'][:, DETECTOR_SOLAR_EXCL_IDX] = np.deg2rad(30)
    sim_data['satellites']['detector'][:, DETECTOR_LUNAR_EXCL_IDX] = np.deg2rad(30)
    sim_data['satellites']['detector'][:, DETECTOR_EARTH_EXCL_IDX] = np.deg2rad(10)

    print("Calculating celestial body positions...")
    sim_data = celestial_update(sim_data, sim_start_time)

    print("Generating exclusion table (this may take a moment)...")
    original_fixed_points = sim_data['fixedpoints']['position']
    sim_data['fixedpoints']['position'] = original_fixed_points[:200]

    update_visibility_table(sim_data)
    exclusion_matrix = sim_data['fixedpoints']['visibility']

    sim_data['fixedpoints']['position'] = original_fixed_points

    print("Exclusion table generated.")

    print("Displaying results as a heatmap...")
    fig = go.Figure(data=go.Heatmap(
        z=exclusion_matrix.T,
        colorscale=[[0, 'red'], [1, 'green']],
        showscale=False
    ))

    fig.update_layout(
        title='Satellite vs. Fixed Point Exclusion Status (0=Clear, 1=Excluded)',
        xaxis_title='Fixed Point Index',
        yaxis_title='Satellite Index',
        yaxis=dict(autorange='reversed')
    )

    return fig
\end{Verbatim}
\subsection{\texttt{demo\_fixedpoints.py}}
\begin{Verbatim}
from datetime import datetime, timezone
import plotly.graph_objects as go

from simulation import create_empty_simulation, add_fixed_points
from plotting_3d import plot_3d_scatter

def demo_fixedpoints() -> go.Figure:
    """
    Demonstrates the fixedpoints data structure by plotting it in 3D.

    This function generates and returns a Plotly figure object but does not
    display it. The caller is responsible for rendering the plot.

    Returns:
        The Plotly figure object.
    """
    print("\n--- Starting Demo Fixedpoints ---")
    sim_start_time = datetime(2025, 7, 27, 22, 27, 0, tzinfo=timezone.utc)

    sim_data = create_empty_simulation(sim_start_time)
    add_fixed_points(sim_data, 10000)

    fixed_positions = sim_data['fixedpoints']['position']

    print(f"Plotting {len(fixed_positions)} fixed points.")

    fig = plot_3d_scatter(
        positions=fixed_positions,
        title="Fixed Points Distribution",
        plot_time=sim_start_time,
        labels=[f"Point {i}" for i in range(len(fixed_positions))],
        marker_size=1,
        trace_name='Fixed Points'
    )
    return fig
\end{Verbatim}
\subsection{\texttt{demo\_lambertian.py}}
\begin{Verbatim}
import numpy as np
import plotly.graph_objects as go
from lambertian import lambertiansphere

def demo_lambertian():
    """
    Runs a demonstration of the lambertiansphere function,
    including example calculations and a plot.
    """
    SATELLITE_ALBEDO = 0.2
    SATELLITE_RADIUS = 1.5

    print(f"--- Simulating a sphere with Albedo={SATELLITE_ALBEDO} and Radius={SATELLITE_RADIUS}m ---\n")

    print("--- Example 1: Full Illumination ---")
    vec_sun_1 = np.array([1, 0, 0])
    vec_obs_1 = np.array([1, 0, 0])
    brightness_1 = lambertiansphere(vec_sun_1, vec_obs_1, SATELLITE_ALBEDO, SATELLITE_RADIUS)
    angle_1 = np.rad2deg(np.arccos(np.dot(vec_sun_1, vec_obs_1)))
    print(f"Phase Angle: {angle_1:.2f} degrees")
    print(f"Effective Brightness: {brightness_1:.4f} m^2\n")

    print("--- Example 2: Half Illumination ---")
    vec_sun_2 = np.array([1, 0, 0])
    vec_obs_2 = np.array([0, 1, 0])
    brightness_2 = lambertiansphere(vec_sun_2, vec_obs_2, SATELLITE_ALBEDO, SATELLITE_RADIUS)
    angle_2 = np.rad2deg(np.arccos(np.dot(vec_sun_2/np.linalg.norm(vec_sun_2),
vec_obs_2/np.linalg.norm(vec_obs_2))))
    print(f"Phase Angle: {angle_2:.2f} degrees")
    print(f"Effective Brightness: {brightness_2:.4f} m^2\n")

    print("--- Example 3: No Illumination ---")
    vec_sun_3 = np.array([1, 0, 0])
    vec_obs_3 = np.array([-1, 0, 0])
    brightness_3 = lambertiansphere(vec_sun_3, vec_obs_3, SATELLITE_ALBEDO, SATELLITE_RADIUS)
    angle_3 = np.rad2deg(np.arccos(np.dot(vec_sun_3, vec_obs_3)))
    print(f"Phase Angle: {angle_3:.2f} degrees")
    print(f"Effective Brightness: {brightness_3:.4f} m^2\n")

    print("\n--- Generating Plot Data ---")
    angles_deg = np.linspace(0, 180, 200)
    angles_rad = np.deg2rad(angles_deg)
    brightness_values = []
    plot_vec_light = np.array([1, 0, 0])
    for angle in angles_rad:
        plot_vec_obs = np.array([np.cos(angle), np.sin(angle), 0])
        brightness = lambertiansphere(plot_vec_light, plot_vec_obs, SATELLITE_ALBEDO, SATELLITE_RADIUS)
        brightness_values.append(brightness)

    fig = go.Figure()
    fig.add_trace(go.Scatter(x=angles_deg, y=brightness_values, mode='lines', name='Effective Brightness'))
    title_text = f'Lambertian Sphere Brightness vs. Phase Angle<br><sup>Albedo={SATELLITE_ALBEDO}, Radius={SATELLITE_RADIUS}m</sup>'
    fig.update_layout(
        title_text=title_text,
        xaxis_title="Physical Phase Angle (degrees)",
        yaxis_title="Effective Brightness (m^2)",
        template="plotly_white",
        xaxis=dict(range=[0, 180]),
        yaxis=dict(range=[0, max(brightness_values) * 1.1])
    )

    print("\n--- Returning Plot ---")
    return fig
\end{Verbatim}
\subsection{\texttt{demo\_pointing\_plot.py}}
\begin{Verbatim}
from datetime import datetime, timezone
import plotly.graph_objects as go

from demo_common import initialize_standard_simulation
from plotting_vectors import plot_pointing_vectors

def demo_pointing_plot() -> go.Figure:
    """
    Demonstrates the plot_pointing_vectors function.

    This function generates and returns a Plotly figure object but does not
    display it. The caller is responsible for rendering the plot.

    Returns:
        The Plotly figure object.
    """
    print("\n--- Starting Demo: Pointing Vector Plot ---")
    sim_start_time = datetime(2025, 8, 1, 12, 0, 0, tzinfo=timezone.utc)
    sim_data = initialize_standard_simulation(sim_start_time)

    fig = plot_pointing_vectors(
        data_struct=sim_data,
        title="Satellite Positions with Radially Outward Pointing Vectors",
        plot_time=sim_start_time
    )
    return fig
\end{Verbatim}
\subsection{\texttt{demo\_pointing\_sequence.py}}
\begin{Verbatim}
import numpy as np
from datetime import datetime, timezone, timedelta
import plotly.graph_objects as go

from simulation import create_empty_simulation
from observatories import add_observatories
from propagation import add_satellites_from_tle
from constants import POINTING_COUNT_IDX, POINTING_PLACE_IDX
from pointing import generate_pointing_sphere, update_satellite_pointing, pointing_place_update
from plotting_vectors import plot_pointing_vectors

def demo_pointing_sequence() -> go.Figure:
    """
    Demonstrates the satellite pointing sequence functionality.
    """
    print("\n--- Starting Demo: Pointing Sequence ---")
    sim_start_time = datetime(2025, 8, 1, 12, 0, 0, tzinfo=timezone.utc)

    # Initialize a simulation with 3 satellites
    sim_data = create_empty_simulation(sim_start_time)
    
    # Create a dummy TLE file for 3 satellites
    tle_data = """SAT-1
1 90401U 25007A   25210.50000000  .00000000  00000-0  00000-0 0  9991
2 90401   0.0500  45.0000 0001000  90.0000  20.0000  1.00270000    11
SAT-2
1 90402U 25007B   25210.50000000  .00000000  00000-0  00000-0 0  9991
2 90402   0.0500  45.0000 0001000  90.0000  20.0000  1.00270000    11
SAT-3
1 90403U 25007C   25210.50000000  .00000000  00000-0  00000-0 0  9991
2 90403   0.0500  45.0000 0001000  90.0000  20.0000  1.00270000    11
"""
    dummy_tle_path = "dummy_tle_pointing.txt"
    with open(dummy_tle_path, "w") as f:
        f.write(tle_data)

    add_satellites_from_tle(sim_data, dummy_tle_path, 'satellites')

    # Generate pointing spheres
    generate_pointing_sphere(sim_data, 10)
    generate_pointing_sphere(sim_data, 20)

    # Assign pointing counts to satellites
    pointing_state = sim_data['satellites']['pointing_state']
    pointing_state[0, POINTING_COUNT_IDX] = 10
    pointing_state[1, POINTING_COUNT_IDX] = 20
    # Satellite 2 will have a pointing_count of 0 and should not move

    print("Initial pointing vectors:")
    update_satellite_pointing(sim_data)
    print(sim_data['satellites']['pointing'])

    # --- Create a figure to animate ---
    fig = go.Figure()

    # Initial plot
    vectors = sim_data['satellites']['pointing']
    fig.add_trace(go.Scatter3d(
        x=vectors[:, 0], y=vectors[:, 1], z=vectors[:, 2],
        mode='markers', marker=dict(size=10, color=['red', 'green', 'blue']),
        name='T=0'
    ))

    # Simulation loop
    for t in range(1, 5):
        print(f"\n--- Time Step {t} ---")
        print("Pointing state:")
        print(sim_data['satellites']['pointing_state'])
        pointing_place_update(sim_data)
        update_satellite_pointing(sim_data)
        print("Pointing vectors:")
        print(sim_data['satellites']['pointing'])

        vectors = sim_data['satellites']['pointing']
        fig.add_trace(go.Scatter3d(
            x=vectors[:, 0], y=vectors[:, 1], z=vectors[:, 2],
            mode='markers', marker=dict(size=10, color=['red', 'green', 'blue'], opacity=0.5),
            name=f'T={t}'
        ))

    # Add a unit sphere for reference
    u = np.linspace(0, 2 * np.pi, 100)
    v = np.linspace(0, np.pi, 100)
    x = np.outer(np.cos(u), np.sin(v))
    y = np.outer(np.sin(u), np.sin(v))
    z = np.outer(np.ones(np.size(u)), np.cos(v))
    fig.add_trace(go.Surface(x=x, y=y, z=z, colorscale='Blues', showscale=False, opacity=0.1))

    fig.update_layout(
        title="Satellite Pointing Sequence",
        scene=dict(
            xaxis_title='X', yaxis_title='Y', zaxis_title='Z',
            aspectmode='data'
        )
    )

    return fig

if __name__ == '__main__':
    fig = demo_pointing_sequence()
    fig.show()
\end{Verbatim}
\subsection{\texttt{demo\_pointing\_vectors.py}}
\begin{Verbatim}
import plotly.graph_objects as go
from pointing_vectors import pointing_vectors, plot_vectors_on_sphere

def demo_pointing_vectors() -> go.Figure:
    """
    Demonstrates the generation and plotting of pointing vectors.

    Returns:
        A Plotly figure object.
    """
    print("--- Running Pointing Vectors Demo ---")
    num_vectors = 1000
    print(f"Generating {num_vectors} pointing vectors...")
    vectors = pointing_vectors(num_vectors)
    print("Plotting vectors on a sphere...")
    fig = plot_vectors_on_sphere(vectors, f"{num_vectors} Pointing Vectors")
    return fig

if __name__ == '__main__':
    fig = demo_pointing_vectors()
    fig.show()
\end{Verbatim}
\subsection{\texttt{demo\_sky\_scan.py}}
\begin{Verbatim}
import numpy as np
from datetime import datetime, timezone
import plotly.graph_objects as go

from simulation import create_empty_simulation, add_celestial_bodies, DETECTOR_SOLAR_EXCL_IDX,
DETECTOR_LUNAR_EXCL_IDX, DETECTOR_EARTH_EXCL_IDX
from propagation import add_satellites_from_tle, propagate_satellites_new, celestial_update
from visibility import exclusion

def demo_sky_scan() -> go.Figure:
    """
    Performs a sky scan from a GEO satellite to map celestial exclusion zones.
    """
    print("\n--- Starting Demo: Sky Scan from GEO ---")
    sim_start_time = datetime(2025, 8, 1, 12, 0, 0, tzinfo=timezone.utc)

    # a. Initialize a simulation with a single GEO satellite
    tle_data = """GEO-SCAN
1 90401U 25007A   25210.50000000  .00000000  00000-0  00000-0 0  9991
2 90401   0.0500  45.0000 0001000  90.0000  20.0000  1.00270000    11
"""
    dummy_tle_path = "dummy_tle_skyscan.txt"
    with open(dummy_tle_path, "w") as f:
        f.write(tle_data)

    sim_data = create_empty_simulation(sim_start_time)
    add_satellites_from_tle(sim_data, dummy_tle_path, 'satellites')
    add_celestial_bodies(sim_data)

    # Set some reasonable exclusion angles
    sim_data['satellites']['detector'][:, DETECTOR_SOLAR_EXCL_IDX] = np.deg2rad(30)
    sim_data['satellites']['detector'][:, DETECTOR_LUNAR_EXCL_IDX] = np.deg2rad(15)
    sim_data['satellites']['detector'][:, DETECTOR_EARTH_EXCL_IDX] = np.deg2rad(10)

    # b. Update celestial body positions
    sim_data = propagate_satellites_new(sim_data, sim_start_time)
    sim_data = celestial_update(sim_data, sim_start_time)
    print("Initialized GEO satellite and celestial bodies.")

    # c. Create a 2D grid of pointing vectors
    declinations = np.linspace(-np.pi / 2, np.pi / 2, 50)
    right_ascensions = np.linspace(0, 2 * np.pi, 100)

    # d. & e. Iterate and store results
    sky_map = np.zeros((len(declinations), len(right_ascensions)))

    print("Scanning the sky for exclusion zones...")
    for i, dec in enumerate(declinations):
        for j, ra in enumerate(right_ascensions):
            # Convert spherical to Cartesian coordinates for the pointing vector
            x = np.cos(dec) * np.cos(ra)
            y = np.cos(dec) * np.sin(ra)
            z = np.sin(dec)
            pointing_vector = np.array([x, y, z])

            # Update the satellite's pointing vector
            sim_data['satellites']['pointing'][0] = pointing_vector

            # Call the exclusion function
            is_clear = exclusion(sim_data, satellite_index=0)
            sky_map[i, j] = is_clear

    print("Sky scan complete.")

    # f. Use Plotly to create a heatmap
    fig = go.Figure(data=go.Heatmap(
        z=sky_map,
        x=np.rad2deg(right_ascensions),
        y=np.rad2deg(declinations),
        colorscale=[[0, 'red'], [1, 'lightgreen']],
        showscale=False,
        colorbar=dict(
            title='Visibility',
            tickvals=[0, 1],
            ticktext=['Excluded', 'Clear']
        )
    ))

    fig.update_layout(
        title='Sky Exclusion Map from a GEO Satellite',
        xaxis_title='Right Ascension (degrees)',
        yaxis_title='Declination (degrees)',
        yaxis=dict(scaleanchor="x", scaleratio=1)
    )

    return fig

if __name__ == '__main__':
    fig = demo_sky_scan()
    fig.show()
\end{Verbatim}
\subsection{\texttt{generate\_log\_spherical\_points.py}}
\begin{Verbatim}
import numpy as np
from datetime import datetime, timezone
from plotting_3d import plot_3d_scatter


def generate_log_spherical_points(
    num_points: int,
    inner_radius: float,
    outer_radius: float,
    object_size_m: float = 1.0,
    seed: int = None
) -> tuple[np.ndarray, np.ndarray]:
    """
    Generates 3D points with logarithmic radial and uniform angular distribution.

    This function creates a point cloud where point distances from the origin are
    logarithmically spaced. On any given spherical shell, points are distributed
    uniformly using the Fibonacci lattice method. Each point is associated with a
    specified object size.

    Args:
        num_points: The total number of points to generate.
        inner_radius: The minimum distance from the origin (must be positive).
        outer_radius: The maximum distance from the origin (must be >= inner_radius).
        object_size_m: The size in meters to be associated with each point.
                       Defaults to 1.0.
        seed: An optional integer to seed the random number generator for
              reproducible shuffling.

    Returns:
        A tuple containing:
        - A NumPy array of shape (num_points, 3) for the Cartesian coordinates.
        - A NumPy array of shape (num_points,) for the object size in meters.
    """
    # Input validation
    if not isinstance(num_points, int) or num_points <= 0:
        raise ValueError("num_points must be a positive integer.")
    if not (isinstance(inner_radius, (int, float)) and inner_radius > 0):
        raise ValueError("inner_radius must be a positive number.")
    if not (isinstance(outer_radius, (int, float)) and outer_radius >= inner_radius):
        raise ValueError("outer_radius must be a positive number and >= inner_radius.")

    # --- 1. Generate uniformly distributed unit vectors (Fibonacci Lattice) ---
    # Create an array of indices for each point
    indices = np.arange(0, num_points, dtype=float) + 0.5

    # Uniformly distribute points along the z-axis (cos(phi))
    z = 1 - 2 * indices / num_points

    # Calculate the radius in the xy-plane for each point
    radius_xy = np.sqrt(1 - z**2)

    # Calculate the azimuthal angle using the golden angle
    golden_angle = np.pi * (3. - np.sqrt(5.))
    theta = golden_angle * indices

    # Convert to Cartesian coordinates for the unit sphere
    x = radius_xy * np.cos(theta)
    y = radius_xy * np.sin(theta)

    # Stack into an (N, 3) array of unit vectors
    unit_vectors = np.stack([x, y, z], axis=1)

    # --- 2. Generate logarithmically spaced radii ---
    # The start and stop parameters for logspace are exponents
    start_exp = np.log10(inner_radius)
    stop_exp = np.log10(outer_radius)
    radii = np.logspace(start_exp, stop_exp, num_points)

    # --- 3. Shuffle radii to decouple radius from spherical position ---
    # This prevents the smallest radii from always mapping to points near the
    # north pole, creating a more uniform visual distribution of densities.
    rng = np.random.default_rng(seed)
    rng.shuffle(radii)

    # --- 4. Scale unit vectors by radii using broadcasting ---
    # Reshape radii to (N, 1) to broadcast with (N, 3) unit_vectors
    points = unit_vectors * radii[:, np.newaxis]

    # --- 5. Create an array for the object sizes ---
    sizes = np.full(num_points, object_size_m, dtype=float)

    return points, sizes

if __name__ == '__main__':
    # --- Demo of Point Generation and Visualization ---
    # This provides a consistent demonstration entry point, matching vibevolts_demo.py
    NUM_POINTS = 100
    # Using the same radii as the main simulation for consistency
    INNER_RADIUS = 2000000
    OUTER_RADIUS = 84328000
    OBJECT_SIZE = 1.0

    print(f"Generating {NUM_POINTS} points from radius {INNER_RADIUS} to {OUTER_RADIUS}...")
    generated_points, _ = generate_log_spherical_points(
        num_points=NUM_POINTS,
        inner_radius=INNER_RADIUS,
        outer_radius=OUTER_RADIUS,
        object_size_m=OBJECT_SIZE,
        seed=42
    )
    print("Point generation complete.\n")

    # Use the unified visualization function
    plot_3d_scatter(
        positions=generated_points,
        title="Fixed Points Distribution (from generate_log_spherical_points.py)",
        plot_time=datetime.now(timezone.utc),
        labels=[f"Point {i}" for i in range(len(generated_points))],
        marker_size=1,
        trace_name='Fixed Points'
    )
\end{Verbatim}
\subsection{\texttt{generate\_report.py}}
\begin{Verbatim}
import plotly.graph_objects as go

from vibevolts_demo import (
    demo1,
    demo2,
    demo3,
    demo4,
    demo_fixedpoints,
    demo_exclusion_table,
    demo_pointing_plot,
)

def generate_demo_html_report():
    """
    Runs all plotting demos and saves the output to a single HTML file.
    """
    print("Generating HTML report for all demos...")

    # A list of all demo functions that produce plots
    plotting_demos = [
        ("Demo 1: Satellite Positions", demo1),
        ("Demo 2: Satellite and Celestial Positions", demo2),
        ("Demo 3: Single LEO Satellite Trajectory", demo3),
        ("Demo 4: Single GEO Satellite Trajectory", demo4),
        ("Demo Fixed Points", demo_fixedpoints),
        ("Demo Exclusion Table", demo_exclusion_table),
        ("Demo Pointing Plot", demo_pointing_plot),
    ]

    # Collect all the figure objects
    figures = []
    for name, demo_func in plotting_demos:
        print(f"Running: {name}")
        fig = demo_func(show_plot=False)
        # Add a title to the figure layout if it doesn't have one
        if fig.layout.title.text is None or fig.layout.title.text == '':
            fig.update_layout(title_text=name)
        figures.append(fig)

    # Write all figures to a single HTML file
    report_filename = "demo_plots.html"
    print(f"Writing report to {report_filename}...")
    with open(report_filename, 'w') as f:
        f.write("<html><head><title>Vibevolts Demo Report</title></head><body>")
        f.write("<h1>Vibevolts Demo Report</h1>")
        # Include plotly.js from CDN
        f.write("<script src='https://cdn.plot.ly/plotly-latest.min.js'></script>")

        for i, fig in enumerate(figures):
            # Convert figure to an HTML div
            fig_html = fig.to_html(full_html=False, include_plotlyjs=False)
            f.write(f"<h2>{fig.layout.title.text}</h2>")
            f.write(fig_html)
            if i < len(figures) - 1:
                f.write("<hr>") # Add a horizontal rule between plots

        f.write("</body></html>")

    print("Report generation complete.")

if __name__ == '__main__':
    generate_demo_html_report()
\end{Verbatim}
\subsection{\texttt{lambertian.py}}
\begin{Verbatim}
import numpy as np

def lambertiansphere(
    vec_from_sphere_to_light: np.ndarray,
    vec_from_sphere_to_observer: np.ndarray,
    albedo: float,
    radius: float
) -> float:
    """
    Calculates the effective brightness of a
    Lambertian sphere.

    This function determines the apparent brightness of
    a diffusely reflecting sphere based on the angle
    between the light source and the observer, the
    sphere's albedo (reflectivity), and its size.

    Args:
        vec_from_sphere_to_light: A 3-element NumPy
            array representing the direction vector from
            the sphere to the light source.
        vec_from_sphere_to_observer: A 3-element NumPy
            array representing the direction vector from
            the sphere to the observer.
        albedo: The fraction of incident light that is
            reflected (0.0 to 1.0).
        radius: The radius of the sphere in meters.

    Returns:
        The effective brightness cross-section in
        square meters. This value is proportional to
        the total light reflected towards the observer.
    """
    if not 0.0 <= albedo <= 1.0:
        raise ValueError("Albedo must be between 0.0 and 1.0.")
    if radius < 0:
        raise ValueError("Radius cannot be negative.")

    norm_light = np.linalg.norm(vec_from_sphere_to_light)
    norm_observer = np.linalg.norm(vec_from_sphere_to_observer)

    if norm_light == 0 or norm_observer == 0:
        raise ValueError("Input vectors cannot have zero length.")

    unit_vec_light = vec_from_sphere_to_light / norm_light
    unit_vec_observer = vec_from_sphere_to_observer / norm_observer

    cos_alpha = np.dot(unit_vec_light, unit_vec_observer)
    cos_alpha = np.clip(cos_alpha, -1.0, 1.0)
    alpha = np.arccos(cos_alpha)

    term1 = np.sin(alpha)
    term2 = (np.pi - alpha) * np.cos(alpha)
    phase_function_value = (2 / (3 * np.pi)) * (term1 + term2)

    cross_sectional_area = np.pi * (radius ** 2)

    effective_brightness = (
        albedo *
        cross_sectional_area *
        phase_function_value
    )

    return effective_brightness
\end{Verbatim}
\subsection{\texttt{observatories.py}}
\begin{Verbatim}
import numpy as np
from typing import Dict, Any

def add_observatories(sim_data: Dict[str, Any], num_observatories: int) -> None:
    """
    Adds observatory data structures to the simulation data.

    Args:
        sim_data: The main simulation data dictionary.
        num_observatories: The number of observatories to add.
    """
    if not isinstance(num_observatories, int) or num_observatories < 0:
        raise ValueError("num_observatories must be a non-negative integer.")

    sim_data['counts']['observatories'] = num_observatories
    sim_data['observatories'] = {
        'position': np.zeros((num_observatories, 3), dtype=float),
        'velocity': np.zeros((num_observatories, 3), dtype=float),
        'acceleration': np.zeros((num_observatories, 3), dtype=float),
        'pointing': np.zeros((num_observatories, 3), dtype=float),
        'detector': np.zeros((num_observatories, 7), dtype=float),
    }
\end{Verbatim}
\subsection{\texttt{plotting\_3d.py}}
\begin{Verbatim}
import numpy as np
from datetime import datetime
from typing import List, Optional

import plotly.graph_objects as go
from astropy.time import Time
from astropy.coordinates import GCRS, EarthLocation
import astropy.units as u

EARTH_RADIUS = 6378137.0

def plot_3d_scatter(
    positions: np.ndarray,
    title: str,
    plot_time: datetime,
    labels: Optional[List[str]] = None,
    marker_size: int = 1,
    trace_name: str = 'Points'
) -> go.Figure:
    """
    Creates a 3D plot of object positions.
    """
    if positions.ndim != 2 or positions.shape[1] != 3:
        raise ValueError("positions array must have shape (n, 3).")

    fig = go.Figure()

    fig.add_trace(go.Scatter3d(
        x=positions[:, 0], y=positions[:, 1], z=positions[:, 2],
        mode='markers',
        marker=dict(
            size=marker_size,
            color=np.arange(len(positions)),
            colorscale='Viridis',
            opacity=0.8
        ),
        text=labels,
        hoverinfo='text' if labels else 'none',
        name=trace_name
    ))

    u_sphere = np.linspace(0, 2 * np.pi, 100)
    v_sphere = np.linspace(0, np.pi, 100)
    x_earth = EARTH_RADIUS * np.outer(np.cos(u_sphere), np.sin(v_sphere))
    y_earth = EARTH_RADIUS * np.outer(np.sin(u_sphere), np.sin(v_sphere))
    z_earth = EARTH_RADIUS * np.outer(np.ones(np.size(u_sphere)), np.cos(v_sphere))
    fig.add_trace(go.Surface(
        x=x_earth, y=y_earth, z=z_earth,
        colorscale='Blues', showscale=False, opacity=0.5, name='Earth'
    ))

    theta = np.linspace(0, 2 * np.pi, 100)
    x_eq = EARTH_RADIUS * np.cos(theta)
    y_eq = EARTH_RADIUS * np.sin(theta)
    z_eq = np.zeros_like(theta)
    fig.add_trace(go.Scatter3d(
        x=x_eq, y=y_eq, z=z_eq, mode='lines',
        line=dict(color='green', width=3), name='Equator'
    ))

    fig.add_trace(go.Scatter3d(
        x=[0], y=[0], z=[EARTH_RADIUS * 1.1], mode='text',
        text=['N'], textfont=dict(size=15, color='red'), name='North Pole'
    ))

    lat_es, lon_es = 33.92 * u.deg, -118.42 * u.deg
    el_segundo_loc = EarthLocation.from_geodetic(lon=lon_es, lat=lat_es)
    itrs_coords = el_segundo_loc.get_itrs(obstime=Time(plot_time))
    gcrs_coords = itrs_coords.transform_to(GCRS(obstime=Time(plot_time)))
    es_pos = gcrs_coords.cartesian.xyz.to(u.m).value * 1.05
    fig.add_trace(go.Scatter3d(
        x=[es_pos[0]], y=[es_pos[1]], z=[es_pos[2]], mode='text',
        text=['ES'], textfont=dict(size=15, color='yellow'), name='El Segundo'
    ))

    fig.update_layout(
        title=title,
        scene=dict(
            xaxis_title='X (m)', yaxis_title='Y (m)', zaxis_title='Z (m)',
            aspectmode='data'
        ),
        margin=dict(r=20, b=10, l=10, t=40),
        legend_title_text='Objects'
    )

    return fig
\end{Verbatim}
\subsection{\texttt{plotting\_vectors.py}}
\begin{Verbatim}
import numpy as np
from datetime import datetime
from typing import Dict, Any

import plotly.graph_objects as go

EARTH_RADIUS = 6378137.0

def plot_pointing_vectors(
    data_struct: Dict[str, Any],
    title: str,
    plot_time: datetime
) -> go.Figure:
    """
    Creates a 3D plot of satellites with pointing vectors.
    """
    sat_positions = data_struct['satellites']['position'].copy()
    sat_pointing = data_struct['satellites']['pointing'].copy()
    num_sats = data_struct['counts']['satellites']

    if num_sats == 0:
        print("No satellites to plot.")
        return go.Figure()

    fig = go.Figure()

    fig.add_trace(go.Scatter3d(
        x=sat_positions[:, 0], y=sat_positions[:, 1], z=sat_positions[:, 2],
        mode='markers',
        marker=dict(size=5, color='blue', opacity=0.8),
        name='Satellites'
    ))

    vector_scale = 0.5 * EARTH_RADIUS
    for i in range(num_sats):
        start_point = sat_positions[i]
        pointing_vec = sat_pointing[i]
        norm = np.linalg.norm(pointing_vec)
        unit_vec = pointing_vec / norm if norm > 0 else np.array([0, 0, 0])
        end_point = start_point + unit_vec * vector_scale
        fig.add_trace(go.Scatter3d(
            x=[start_point[0], end_point[0]],
            y=[start_point[1], end_point[1]],
            z=[start_point[2], end_point[2]],
            mode='lines',
            line=dict(color='red', width=3),
            showlegend=False
        ))

    u_sphere = np.linspace(0, 2 * np.pi, 100)
    v_sphere = np.linspace(0, np.pi, 100)
    x_earth = EARTH_RADIUS * np.outer(np.cos(u_sphere), np.sin(v_sphere))
    y_earth = EARTH_RADIUS * np.outer(np.sin(u_sphere), np.sin(v_sphere))
    z_earth = EARTH_RADIUS * np.outer(np.ones(np.size(u_sphere)), np.cos(v_sphere))
    fig.add_trace(go.Surface(
        x=x_earth, y=y_earth, z=z_earth,
        colorscale='Blues', showscale=False, opacity=0.5, name='Earth'
    ))

    fig.update_layout(
        title=title,
        scene=dict(
            xaxis_title='X (m)', yaxis_title='Y (m)', zaxis_title='Z (m)',
            aspectmode='data'
        ),
        margin=dict(r=20, b=10, l=10, t=40)
    )

    return fig
\end{Verbatim}
\subsection{\texttt{pointing.py}}
\begin{Verbatim}
import numpy as np
from typing import Dict, Any

from constants import POINTING_COUNT_IDX, POINTING_PLACE_IDX
from pointing_vectors import pointing_vectors

def pointing_place_update(data_struct: Dict[str, Any]) -> None:
    """
    Increments the pointing place for all satellites, wrapping around if necessary.
    Basically this usges data_struct['satellites']['pointing_state'] and wraps to
    0 if its     POINTING_COUNT_IDX or greater.
    """

    pointing_state = data_struct['satellites']['pointing_state']
    pointing_counts = pointing_state[:, POINTING_COUNT_IDX]

    # Increment pointing place
    pointing_state[:, POINTING_PLACE_IDX] += 1

    # Wrap around where place >= count
    wrap_around_indices = np.where(pointing_state[:, POINTING_PLACE_IDX] >= pointing_counts)
    pointing_state[wrap_around_indices, POINTING_PLACE_IDX] = 0

def jerk(data_struct: Dict[str, Any], satellite_number: int) -> Dict[str, Any]:
    """
    Moves the pointing vector of a specific satellite by 0.3 radians in a
    random direction.

    This function applies a random rotation to the satellite's pointing vector
    using a simplified version of Rodrigues' rotation formula.

    Args:
        data_struct: The main simulation data dictionary.
        satellite_number: The index of the satellite to modify.

    Returns:
        The modified data_struct with the updated pointing vector.
    """
    p = data_struct['satellites']['pointing'][satellite_number]
    p_norm = p / np.linalg.norm(p)

    while True:
        r = np.random.rand(3) - 0.5
        k = np.cross(p_norm, r)
        norm_k = np.linalg.norm(k)
        if norm_k > 1e-9:
            k_hat = k / norm_k
            break

    theta = 0.3
    cos_theta = np.cos(theta)
    sin_theta = np.sin(theta)

    p_new = p_norm * cos_theta + np.cross(k_hat, p_norm) * sin_theta

    data_struct['satellites']['pointing'][satellite_number] = p_new / np.linalg.norm(p_new)

    return data_struct

def generate_pointing_sphere(data_struct: Dict[str, Any], n_points: int) -> None:
    """
    Generates a pointing sphere with n_points and stores it in the data_struct.
    The index is ['pointing_spheres'][n]
    If a sphere with the same number of points already exists, this function does nothing.
    """
    if n_points not in data_struct['pointing_spheres']:
        print(f"Generating pointing sphere with {n_points} points...")
        data_struct['pointing_spheres'][n_points] = pointing_vectors(n_points)

def update_satellite_pointing(data_struct: Dict[str, Any]) -> None:
    """
    Updates the pointing vector for each satellite based on its pointing state.
    """
    num_sats = data_struct['counts']['satellites']
    if num_sats == 0:
        return

    pointing_state = data_struct['satellites']['pointing_state']
    pointing_vectors_all = data_struct['satellites']['pointing']

    for i in range(num_sats):
        count = int(pointing_state[i, POINTING_COUNT_IDX])
        place = int(pointing_state[i, POINTING_PLACE_IDX])

        if count > 0:
            if count not in data_struct['pointing_spheres']:
                raise ValueError(f"Pointing sphere for {count} points not generated.")

            grid = data_struct['pointing_spheres'][count]
            pointing_vectors_all[i] = grid[place]

def find_and_jerk_blind_satellites(data_struct: Dict[str, Any]) -> Dict[str, Any]:
    """
    Finds satellites with no visibility and applies the 'jerk' function to them.

    This function iterates through the visibility table. If any satellite (column)
    has no visible fixed points (i.e., the column sum is 0), the `jerk`
    function is called to randomly adjust its pointing vector.

    Args:
        data_struct: The main simulation data dictionary.

    Returns:
        The modified data_struct.
    """
    visibility_table = data_struct['fixedpoints']['visibility']

    column_sums = np.sum(visibility_table, axis=0)
    blind_satellite_indices = np.where(column_sums == 0)[0]

    for sat_index in blind_satellite_indices:
        print(f"Satellite {sat_index} has no visible points. Applying jerk.")
        data_struct = jerk(data_struct, sat_index)

    return data_struct
\end{Verbatim}
\subsection{\texttt{pointing\_vectors.py}}
\begin{Verbatim}
import numpy as np
import plotly.graph_objects as go


def pointing_vectors(n: int) -> np.ndarray:
    """
    Generates n equally spaced points on a unit sphere using the Fibonacci lattice algorithm.

    Args:
        n: The number of points to generate.

    Returns:
        A NumPy array of shape (n, 3) for the Cartesian coordinates of the points.
    """
    if not isinstance(n, int) or n <= 0:
        raise ValueError("n must be a positive integer.")

    indices = np.arange(0, n, dtype=float) + 0.5
    z = 1 - 2 * indices / n
    radius_xy = np.sqrt(1 - z**2)
    golden_angle = np.pi * (3. - np.sqrt(5.))
    theta = golden_angle * indices
    x = radius_xy * np.cos(theta)
    y = radius_xy * np.sin(theta)

    unit_vectors = np.stack([x, y, z], axis=1)
    return unit_vectors

def plot_vectors_on_sphere(vectors: np.ndarray, title: str) -> go.Figure:
    """
    Creates a 3D plot of vectors on a sphere.

    Args:
        vectors: A NumPy array of shape (n, 3) representing the vectors.
        title: The title of the plot.

    Returns:
        A Plotly figure object.
    """
    if vectors.ndim != 2 or vectors.shape[1] != 3:
        raise ValueError("vectors array must have shape (n, 3).")

    fig = go.Figure()

    fig.add_trace(go.Scatter3d(
        x=vectors[:, 0], y=vectors[:, 1], z=vectors[:, 2],
        mode='markers',
        marker=dict(
            size=2,
            color='red',
            opacity=0.8
        ),
        name='Vectors'
    ))

    u_sphere = np.linspace(0, 2 * np.pi, 100)
    v_sphere = np.linspace(0, np.pi, 100)
    x_earth = 1.0 * np.outer(np.cos(u_sphere), np.sin(v_sphere))
    y_earth = 1.0 * np.outer(np.sin(u_sphere), np.sin(v_sphere))
    z_earth = 1.0 * np.outer(np.ones(np.size(u_sphere)), np.cos(v_sphere))
    fig.add_trace(go.Surface(
        x=x_earth, y=y_earth, z=z_earth,
        colorscale='Blues', showscale=False, opacity=0.5, name='Sphere'
    ))

    fig.update_layout(
        title=title,
        scene=dict(
            xaxis_title='X', yaxis_title='Y', zaxis_title='Z',
            aspectmode='data'
        ),
        margin=dict(r=20, b=10, l=10, t=40),
        legend_title_text='Objects'
    )

    return fig
\end{Verbatim}
\subsection{\texttt{propagation.py}}
\begin{Verbatim}
import numpy as np
from datetime import datetime, timezone
from typing import Dict, Any, List, Tuple
from sgp4.api import Satrec
from astropy.time import Time
from astropy.coordinates import get_body, GCRS
import astropy.units as u

from constants import (
    ORBITAL_A_IDX, ORBITAL_E_IDX, ORBITAL_I_IDX,
    ORBITAL_RAAN_IDX, ORBITAL_ARGP_IDX, ORBITAL_M_IDX
)

def add_satellites_from_tle(sim_data: Dict[str, Any], tle_file_path: str, sat_category: str) -> None:
    """
    Adds and initializes a category of satellites from a TLE file.
    Mote although the TLEs are loaded, positions etc. are not.

    Args:
        sim_data: The main simulation data dictionary.
        tle_file_path: Path to the TLE file.
        sat_category: The key for this satellite category (e.g., 'satellites').

    Data added to the set_category element of sim_data
    position
    velocity
    acceleration
    orbital_elements
    epochs
    pointing
    """
    orbital_elements, epochs = readtle(tle_file_path)
    num_sats = len(epochs)

    sim_data['counts'][sat_category] = num_sats
    sim_data[sat_category] = {
        'position': np.zeros((num_sats, 3), dtype=float),
        'velocity': np.zeros((num_sats, 3), dtype=float),
        'acceleration': np.zeros((num_sats, 3), dtype=float),
        'orbital_elements': orbital_elements,
        'epochs': epochs,
        'pointing': np.zeros((num_sats, 3), dtype=float),
        'pointing_state': np.zeros((num_sats, 2), dtype=int),
        'detector': np.zeros((num_sats, 7), dtype=float),
    }

def celestial_update(data_struct: Dict[str, Any], time_date: datetime) -> Dict[str, Any]:
    """
    Calculates and updates the positions of the Sun and Moon.
    """
    if time_date.tzinfo is None:
        raise ValueError("time_date must be timezone-aware.")

    astro_time = Time(time_date)
    sun_coords = get_body("sun", astro_time)
    sun_gcrs = sun_coords.transform_to(GCRS(obstime=astro_time))
    moon_coords = get_body("moon", astro_time)
    moon_gcrs = moon_coords.transform_to(GCRS(obstime=astro_time))

    celestial_pos = data_struct['celestial']['position']
    celestial_pos[0] = sun_gcrs.cartesian.xyz.to(u.m).value
    celestial_pos[1] = moon_gcrs.cartesian.xyz.to(u.m).value

    return data_struct

def readtle(tle_file_path: str) -> Tuple[np.ndarray, List[datetime]]:
    """
    Reads a TLE file and extracts orbital elements and epochs for each satellite.

    The array returned ahd the orbital elements in "canonical" order.
    """
    orbital_elements_list = []
    epochs_list = []
    with open(tle_file_path, 'r') as f:
        lines = f.readlines()

    for i in range(0, len(lines), 3):
        line1 = lines[i+1].strip()
        line2 = lines[i+2].strip()

        satellite = Satrec.twoline2rv(line1, line2)

        jd, fr = satellite.jdsatepoch, satellite.jdsatepochF
        epoch_dt = Time(jd, fr, format='jd', scale='utc').to_datetime(timezone.utc)
        epochs_list.append(epoch_dt)

        a = satellite.a * satellite.radiusearthkm * 1000.0
        e = satellite.ecco
        inc = satellite.inclo
        raan = satellite.nodeo
        argp = satellite.argpo
        M = satellite.mo

        elements = np.zeros(6)
        elements[ORBITAL_A_IDX] = a
        elements[ORBITAL_E_IDX] = e
        elements[ORBITAL_I_IDX] = inc
        elements[ORBITAL_RAAN_IDX] = raan
        elements[ORBITAL_ARGP_IDX] = argp
        elements[ORBITAL_M_IDX] = M
        orbital_elements_list.append(elements)

    return np.array(orbital_elements_list, dtype=float), epochs_list

def propagate_satellites_new(data_struct: Dict[str, Any], time_date: datetime,
sat_category: str = None) -> Dict[str, Any]:
    """
    Updates satellite positions and pointing vectors based on their orbital elements.
    """
    MU_EARTH = 3.986004418e14
    time_date_timestamp = time_date.timestamp()

    if sat_category:
        categories = [sat_category]
    else:
        categories = ['satellites', 'red_satellites']

    for category in categories:
        if category not in data_struct['counts'] or data_struct['counts'][category] == 0:
            continue

        elements = data_struct[category]['orbital_elements']
        epochs = data_struct[category]['epochs']

        epoch_timestamps = np.array([e.timestamp() for e in epochs])
        delta_t_array = time_date_timestamp - epoch_timestamps

        a = elements[:, ORBITAL_A_IDX]
        e = elements[:, ORBITAL_E_IDX]
        i = elements[:, ORBITAL_I_IDX]
        raan = elements[:, ORBITAL_RAAN_IDX]
        argp = elements[:, ORBITAL_ARGP_IDX]
        M0 = elements[:, ORBITAL_M_IDX]

        n = np.sqrt(MU_EARTH / a**3)
        M = (M0 + n * delta_t_array) % (2 * np.pi)

        E = M.copy()
        for _ in range(10):
            f_E = E - e * np.sin(E) - M
            f_prime_E = 1 - e * np.cos(E)
            f_prime_E[f_prime_E == 0] = 1e-10
            E = E - f_E / f_prime_E

        tan_nu_half = np.sqrt((1 + e) / (1 - e)) * np.tan(E / 2)
        nu = 2 * np.arctan(tan_nu_half)

        r = a * (1 - e * np.cos(E))

        x_pqw = r * np.cos(nu)
        y_pqw = r * np.sin(nu)

        cos_raan = np.cos(raan)
        sin_raan = np.sin(raan)
        cos_argp = np.cos(argp)
        sin_argp = np.sin(argp)
        cos_i = np.cos(i)
        sin_i = np.sin(i)

        P_x = cos_argp * cos_raan - sin_argp * sin_raan * cos_i
        P_y = cos_argp * sin_raan + sin_argp * cos_raan * cos_i
        P_z = sin_argp * sin_i

        Q_x = -sin_argp * cos_raan - cos_argp * sin_raan * cos_i
        Q_y = -sin_argp * sin_raan + cos_argp * cos_raan * cos_i
        Q_z = cos_argp * sin_i

        x_gcrs = x_pqw * P_x + y_pqw * Q_x
        y_gcrs = x_pqw * P_y + y_pqw * Q_y
        z_gcrs = x_pqw * P_z + y_pqw * Q_z

        positions = np.vstack((x_gcrs, y_gcrs, z_gcrs)).T
        data_struct[category]['position'] = positions

        norms = np.linalg.norm(positions, axis=1)[:, np.newaxis]
        norms[norms == 0] = 1.0
        data_struct[category]['pointing'] = positions / norms

    return data_struct
\end{Verbatim}
\subsection{\texttt{radiometry\_calcs.py}}
\begin{Verbatim}
import math
import scipy.integrate
import scipy.constants as const
import numpy as np
import plotly.graph_objects as go

def mag(x: float) -> float:
    """
    Calculates a magnitude value from a linear ratio.
    Uses the formula: magnitude = -2.5 * log10(ratio)
    """
    if x <= 0:
        raise ValueError("Input for mag() must be positive.")
    return -2.5 * math.log10(x)

def amag(x: float) -> float:
    """
    Calculates the linear ratio from a magnitude value.
    This is the inverse of the mag() function.
    Uses the formula: ratio = 10**(-0.4 * magnitude)
    """
    return 10**(-0.4 * x)

def _planck_law(wav_m: float, temp_k: float) -> float:
    """
    Helper function for Planck's law for spectral radiance.
    Args:
        wav_m: Wavelength in meters.
        temp_k: Temperature in Kelvin.
    Returns:
        Spectral radiance in W / (m^2 * sr * m).
    """
    if wav_m <= 0:
        return 0.0
    exponent = (const.h * const.c) / (wav_m * const.k * temp_k)
    if exponent > 700:
        return 0.0
    numerator = 2.0 * const.h * const.c**2
    denominator = (wav_m**5) * (math.expm1(exponent))
    if denominator == 0:
        return 0.0
    return numerator / denominator

def blackbody_flux(
    temperature: float,
    lambda_short: float,
    lambda_long: float
) -> float:
    """
    Numerically computes the integrated spectral radiance of
    a blackbody over a given wavelength band.

    Requires the scipy library.

    Args:
        temperature: The temperature of the blackbody in Kelvin.
        lambda_short: The short wavelength of the band in microns.
        lambda_long: The long wavelength of the band in microns.

    Returns:
        The integrated spectral radiance in units of
        Watts / (m^2 * steradian).
    """
    if lambda_short >= lambda_long:
        raise ValueError("lambda_short must be less than lambda_long.")
    lambda_short_m = lambda_short * 1e-6
    lambda_long_m = lambda_long * 1e-6
    integrated_radiance, _ = scipy.integrate.quad(
        lambda wav: _planck_law(wav, temperature),
        lambda_short_m,
        lambda_long_m
    )
    return integrated_radiance

def stefan_boltzmann_law(temperature: float) -> float:
    """
    Calculates the total power radiated per unit area by a
    blackbody using the Stefan-Boltzmann law.

    Args:
        temperature: The temperature of the blackbody in Kelvin.

    Returns:
        The total radiated power per unit area in W / m^2.
    """
    return const.sigma * (temperature**4)

def plot_blackbody_spectrum(temperature: float):
    """
    Plots the spectral radiance of a blackbody from
    0.5 to 30 microns.

    Args:
        temperature: The temperature of the blackbody in Kelvin.
    """
    wavelengths_microns = np.geomspace(0.5, 30, 500)
    wavelengths_m = wavelengths_microns * 1e-6
    spectral_radiance = [_planck_law(w, temperature) * 1e-6 for w in wavelengths_m]
    fig = go.Figure()
    fig.add_trace(go.Scatter(
        x=wavelengths_microns,
        y=spectral_radiance,
        mode='lines',
        name=f'{temperature} K Blackbody'
    ))
    fig.update_layout(
        title_text=f'Blackbody Spectrum ({temperature} K)',
        xaxis_title="Wavelength (microns)",
        yaxis_title="Spectral Radiance (W / m^2 / sr / micron)",
        xaxis_type="log",
        yaxis_type="log",
        template="plotly_white"
    )
    fig.show()

def plot_blackbody_spectrum_visible_nir(temperature: float):
    """
    Plots the spectral radiance of a blackbody from
    0.1 to 1 micron.

    Args:
        temperature: The temperature of the blackbody in Kelvin.
    """
    wavelengths_microns = np.linspace(0.1, 1.0, 500)
    wavelengths_m = wavelengths_microns * 1e-6
    spectral_radiance = [_planck_law(w, temperature) * 1e-6 for w in wavelengths_m]
    fig = go.Figure()
    fig.add_trace(go.Scatter(
        x=wavelengths_microns,
        y=spectral_radiance,
        mode='lines',
        name=f'{temperature} K Blackbody'
    ))
    fig.update_layout(
        title_text=f'Blackbody Spectrum ({temperature} K) - Visible/NIR',
        xaxis_title="Wavelength (microns)",
        yaxis_title="Spectral Radiance (W / m^2 / sr / micron)",
        template="plotly_white"
    )
    fig.show()
\end{Verbatim}
\subsection{\texttt{radiometry\_data.py}}
\begin{Verbatim}
import scipy.constants as const

# --- Physical Constants ---
AU_M = const.au  # Astronomical Unit in meters
RSUN_M = 6.957e8 # Radius of the Sun in meters

# --- Radiometric Data ---

# Data is compiled for standard astronomical filters.
#
# Magnitudes are apparent magnitudes in the AB or Vega system.
# Sky brightness is for a dark site at new moon, in
# magnitudes per square arcsecond. For space-based IR
# telescopes, this is the zodiacal + telescope background.
# For ground-based IR, it's dominated by thermal emission.
# Wavelengths and bandwidths are in nanometers (nm).
# Zero point is the photon flux for a 0-magnitude
# object in photons per second per square meter.

FILTER_DATA = {
    # Johnson-Cousins UBVRI Filters
    'U': {
        'sun': -26.03,
        'sky': 22.0,
        'space': 22.9,
        'central_wavelength': 365.0,
        'bandwidth': 66.0,
        'zero_point': 4.96e9,
    },
    'B': {
        'sun': -26.13,
        'sky': 22.7,
        'space': 23.3,
        'central_wavelength': 445.0,
        'bandwidth': 94.0,
        'zero_point': 1.36e10,
    },
    'V': {
        'sun': -26.78,
        'sky': 21.9,
        'space': 23,
        'central_wavelength': 551.0,
        'bandwidth': 88.0,
        'zero_point': 8.79e9,
    },
    'R': {
        'sun': -27.11,
        'sky': 21.0,
        'space': 22.0,
        'central_wavelength': 658.0,
        'bandwidth': 138.0,
        'zero_point': 9.74e9,
    },
    'I': {
        'sun': -27.47,
        'sky': 20.0,
        'space': 21.2,
        'central_wavelength': 806.0,
        'bandwidth': 149.0,
        'zero_point': 7.11e9,
    },
    # Near-Infrared JHK Filters
    'J': {
        'sun': -27.91,
        'sky': 16.0,
        'space': 21.5,
        'central_wavelength': 1235.0,
        'bandwidth': 162.0,
        'zero_point': 3.15e9,
    },
    'H': {
        'sun': -28.25,
        'sky': 14.0,
        'space': 20.5,
        'central_wavelength': 1646.0,
        'bandwidth': 251.0,
        'zero_point': 2.35e9,
    },
    'K': {
        'sun': -28.29,
        'sky': 13.0,
        'space': 19.5,
        'central_wavelength': 2159.0,
        'bandwidth': 262.0,
        'zero_point': 1.17e9,
    },
    # SDSS Filters
    'g': {
        'sun': -26.47,
        'sky': 21.8,
        'space': 23.0,
        'central_wavelength': 477.0,
        'bandwidth': 137.9,
        'zero_point': 1.582e10,
    },
    'r': {
        'sun': -26.93,
        'sky': 20.8,
        'space': 22.2,
        'central_wavelength': 623.1,
        'bandwidth': 138.2,
        'zero_point': 1.214e10,
    },
    'i': {
        'sun': -27.05,
        'sky': 20.2,
        'space': 21.5,
        'central_wavelength': 762.5,
        'bandwidth': 153.5,
        'zero_point': 1.103e10,
    },
    'z': {
        'sun': -27.07,
        'sky': 19.0,
        'space': 20.8,
        'central_wavelength': 913.4,
        'bandwidth': 140.9,
        'zero_point': 8.455e9,
    },
    # Ground-Based Mid-IR Filters
    'L': {
        'sun': None,
        'sky': 3.5,
        'central_wavelength': 3500.0,
        'bandwidth': 600.0,
        'zero_point': 9.4e9,
    },
    'M': {
        'sun': None,
        'sky': 0.0,
        'central_wavelength': 4800.0,
        'bandwidth': 600.0,
        'zero_point': 6.8e9,
    },
    'N': {
        'sun': None,
        'sky': -6.0,
        'central_wavelength': 10200.0,
        'bandwidth': 5000.0,
        'zero_point': 2.7e10,
    },
    # JWST MIRI Filters
    'F560W': { 'sun': None, 'sky': 24.4, 'central_wavelength': 5600.0, 'bandwidth': 1000.0,
'zero_point': 9.78e9 },
    'F770W': { 'sun': None, 'sky': 24.4, 'central_wavelength': 7700.0, 'bandwidth': 1950.0,
'zero_point': 1.39e10 },
    'F1000W': { 'sun': None, 'sky': 23.6, 'central_wavelength': 10000.0, 'bandwidth': 1800.0,
'zero_point': 9.86e9 },
    'F1500W': { 'sun': None, 'sky': 21.9, 'central_wavelength': 15000.0, 'bandwidth': 2940.0,
'zero_point': 1.08e10 },
    'F2550W': { 'sun': None, 'sky': 19.4, 'central_wavelength': 25500.0, 'bandwidth': 4050.0,
'zero_point': 8.70e9 }
}
\end{Verbatim}
\subsection{\texttt{show\_geo\_search.py}}
\begin{Verbatim}
import numpy as np
from datetime import datetime, timezone
import plotly.graph_objects as go

from simulation import create_empty_simulation, add_celestial_bodies
from constellation import geos
from propagation import propagate_satellites_new
from pointing import pointing_place_update, update_satellite_pointing
from plotting_vectors import plot_pointing_vectors

def show_geo_search():
    """
    This demo initializes a simulation, adds a GEO constellation, and then
    generates several plots to visualize the satellite pointing updates and
    the RA/Dec history of one satellite.
    """
    sim_start_time = datetime(2025, 7, 27, 22, 27, 0, tzinfo=timezone.utc)
    sim_data = create_empty_simulation(sim_start_time)
    add_celestial_bodies(sim_data)

    geos(sim_data, 12, 0.4)

    propagate_satellites_new(sim_data, sim_start_time)

    update_satellite_pointing(sim_data)
    fig1 = plot_pointing_vectors(sim_data, 'Initial Pointing Vectors', sim_start_time)

    ra_history = []
    dec_history = []

    def record_ra_dec():
        p = sim_data['satellites']['pointing'][0]
        p_norm = p / np.linalg.norm(p)
        ra = np.arctan2(p_norm[1], p_norm[0])
        dec = np.arcsin(p_norm[2])
        ra_history.append(np.rad2deg(ra))
        dec_history.append(np.rad2deg(dec))

    record_ra_dec()

    for _ in range(5):
        pointing_place_update(sim_data)
        update_satellite_pointing(sim_data)
        record_ra_dec()

    fig2 = plot_pointing_vectors(sim_data, 'After 5 Updates', sim_start_time)

    for _ in range(10):
        pointing_place_update(sim_data)
        update_satellite_pointing(sim_data)
        record_ra_dec()

    fig3 = plot_pointing_vectors(sim_data, 'After 15 Updates', sim_start_time)

    for _ in range(85):
        pointing_place_update(sim_data)
        update_satellite_pointing(sim_data)
        record_ra_dec()

    fig4 = go.Figure()
    fig4.add_trace(go.Scatter(y=ra_history, mode='lines', name='RA (deg)'))
    fig4.add_trace(go.Scatter(y=dec_history, mode='lines', name='Dec (deg)'))
    fig4.update_layout(
        title='RA and Dec History of a GEO Satellite',
        xaxis_title='Update Number',
        yaxis_title='Angle (degrees)'
    )

    x = (fig1, fig2, fig3, fig4)
    return fig1, fig2, fig3, fig4
\end{Verbatim}
\subsection{\texttt{simulation.py}}
\begin{Verbatim}
import numpy as np
from datetime import datetime, timezone
from typing import Dict, Any
from generate_log_spherical_points import generate_log_spherical_points
from constants import *

def create_empty_simulation(start_time: datetime, delta_time: float = 60.0) -> Dict[str, Any]:
    """
    Initializes a minimal, empty data structure for a space simulation.

    Args:
        start_time: The starting time and date of the simulation. This must be a
                    timezone-aware datetime object set to UTC.
        delta_time: The time step for the simulation in seconds.

    Returns:
        A dictionary representing the basic simulation state.
    """
    if not isinstance(start_time, datetime):
        raise TypeError("start_time must be a datetime object.")
    if start_time.tzinfo is None:
        raise ValueError("start_time must be timezone-aware. Please set tzinfo.")

    simulation_data: Dict[str, Any] = {
        'start_time': start_time,
        'delta_time': delta_time,
        'counts': {},
        'pointing_spheres': {},
    }
    return simulation_data

def add_celestial_bodies(sim_data: Dict[str, Any]) -> None:
    """
    Adds celestial body structures (for Sun and Moon) to the simulation data.

    Args:
        sim_data: The simulation data dictionary.
    """
    sim_data['counts']['celestial'] = 2  # Sun and Moon
    sim_data['celestial'] = {
        'position': np.zeros((2, 3), dtype=float),
        'velocity': np.zeros((2, 3), dtype=float),
        'acceleration': np.zeros((2, 3), dtype=float),
    }

def add_fixed_points(sim_data: Dict[str, Any], num_points: int = 100) -> None:
    """
    Adds a structure for fixed reference points in the GCRS frame.

    Args:
        sim_data: The simulation data dictionary.
        num_points: The number of fixed points to generate.
    """
    sim_data['counts']['fixedpoints'] = num_points
    sim_data['fixedpoints'] = {
        'position': generate_log_spherical_points(
            num_points=num_points,
            inner_radius=2000000,
            outer_radius=84328000
        )[0],
        'visibility': np.zeros((num_points, 0), dtype=int) # Visibility will be resized later
    }
\end{Verbatim}
\subsection{\texttt{visibility.py}}
\begin{Verbatim}
import numpy as np
from typing import Dict, Any, Tuple, Optional

from constants import (
    EARTH_RADIUS, MOON_RADIUS, DETECTOR_SOLAR_EXCL_IDX,
    DETECTOR_LUNAR_EXCL_IDX, DETECTOR_EARTH_EXCL_IDX
)

def solarexclusion(data_struct: Dict[str, Any]) -> Tuple[np.ndarray, np.ndarray]:
    """
    Calculates solar exclusion for all satellites based on their pointing vectors.

    This function operates in a vectorized manner on all satellites in the
    'satellites' category. It computes the angle between each satellite's
    pointing vector and the vector from the satellite to the Sun.

    Args:
        data_struct: The main simulation data dictionary.

    Returns:
        A tuple containing:
        - exclusion_vector (np.ndarray): An array of the same length as the
          number of satellites. An element is 1 if the satellite is within
          the solar exclusion angle, 0 otherwise.
        - angle_vector (np.ndarray): An array containing the calculated angle
          in radians for each satellite.
    """
    num_sats = data_struct['counts']['satellites']
    if num_sats == 0:
        return np.array([]), np.array([])

    sun_pos = data_struct['celestial']['position'][0]
    sat_pos = data_struct['satellites']['position']
    sat_pointing = data_struct['satellites']['pointing']
    solar_exclusion_angles = data_struct['satellites']['detector'][:, DETECTOR_SOLAR_EXCL_IDX]

    vec_sat_to_sun = sun_pos - sat_pos

    norm_sat_to_sun = np.linalg.norm(vec_sat_to_sun, axis=1)
    norm_sat_pointing = np.linalg.norm(sat_pointing, axis=1)

    valid_norms = (norm_sat_to_sun > 1e-9) & (norm_sat_pointing > 1e-9)

    angle_vector = np.full(num_sats, np.pi)

    if np.any(valid_norms):
        dot_product = np.einsum('ij,ij->i', vec_sat_to_sun[valid_norms], sat_pointing[valid_norms])
        cos_angle = dot_product / (norm_sat_to_sun[valid_norms] * norm_sat_pointing[valid_norms])
        cos_angle = np.clip(cos_angle, -1.0, 1.0)
        angle_vector[valid_norms] = np.arccos(cos_angle)

    exclusion_vector = (angle_vector < solar_exclusion_angles).astype(int)

    return exclusion_vector, angle_vector

def exclusion(
    data_struct: Dict[str, Any],
    satellite_index: int,
    print_debug: bool = False
) -> int:
    """
    Determines if a satellite's pointing vector is excluded by the Sun, Moon, or Earth.

    Args:
        data_struct: The main simulation data dictionary.
        satellite_index: The index of the satellite to check.
        print_debug: If True, prints detailed debug information for the calculation.

    Returns:
        0 if the satellite's view is excluded, 1 otherwise.
    """
    sat_pos = data_struct['satellites']['position'][satellite_index]
    sat_pointing = data_struct['satellites']['pointing'][satellite_index]
    sun_pos = data_struct['celestial']['position'][0]
    moon_pos = data_struct['celestial']['position'][1]

    detector_props = data_struct['satellites']['detector'][satellite_index]
    solar_excl_angle = detector_props[DETECTOR_SOLAR_EXCL_IDX]
    lunar_excl_angle = detector_props[DETECTOR_LUNAR_EXCL_IDX]
    earth_excl_angle = detector_props[DETECTOR_EARTH_EXCL_IDX]

    vec_to_sun = sun_pos - sat_pos
    vec_to_moon = moon_pos - sat_pos
    vec_to_earth = -sat_pos

    dist_to_sun = np.linalg.norm(vec_to_sun)
    dist_to_moon = np.linalg.norm(vec_to_moon)
    dist_to_earth = np.linalg.norm(vec_to_earth)

    u_vec_to_sun = vec_to_sun / dist_to_sun if dist_to_sun > 0 else np.array([0.,0.,0.])
    u_vec_to_moon = vec_to_moon / dist_to_moon if dist_to_moon > 0 else np.array([0.,0.,0.])
    u_vec_to_earth = vec_to_earth / dist_to_earth if dist_to_earth > 0 else np.array([0.,0.,0.])

    norm_pointing = np.linalg.norm(sat_pointing)
    u_sat_pointing = sat_pointing / norm_pointing if norm_pointing > 0 else np.array([0.,0.,0.])

    sun_flag, moon_flag, earth_flag = False, False, False

    cos_angle_sun = np.clip(np.dot(u_sat_pointing, u_vec_to_sun), -1.0, 1.0)
    angle_sun = np.arccos(cos_angle_sun)
    if angle_sun < solar_excl_angle:
        sun_flag = True

    cos_angle_moon = np.clip(np.dot(u_sat_pointing, u_vec_to_moon), -1.0, 1.0)
    angle_moon = np.arccos(cos_angle_moon)
    apparent_radius_moon = np.arctan(MOON_RADIUS / dist_to_moon) if dist_to_moon > 0 else 0
    if (angle_moon - apparent_radius_moon) < lunar_excl_angle:
        moon_flag = True

    cos_angle_earth = np.clip(np.dot(u_sat_pointing, u_vec_to_earth), -1.0, 1.0)
    angle_earth = np.arccos(cos_angle_earth)
    apparent_radius_earth = np.arctan(EARTH_RADIUS / dist_to_earth) if dist_to_earth > 0 else 0
    if (angle_earth - apparent_radius_earth) < earth_excl_angle:
        earth_flag = True

    if print_debug:
        print(f"--- Exclusion Debug for Satellite {satellite_index} ---")
        print(f"  - Sun Flag:   {sun_flag}, Angle: {np.rad2deg(angle_sun):.2f} deg, Excl:
{np.rad2deg(solar_excl_angle):.2f} deg")
        print(f"  - Moon Flag:  {moon_flag}, Angle: {np.rad2deg(angle_moon):.2f} deg, Excl:
{np.rad2deg(lunar_excl_angle):.2f} deg")
        print(f"  - Earth Flag: {earth_flag}, Angle: {np.rad2deg(angle_earth):.2f} deg, Excl:
{np.rad2deg(earth_excl_angle):.2f} deg")

    is_excluded = sun_flag or moon_flag or earth_flag
    return 0 if is_excluded else 1

def update_visibility_table(
    data_struct: Dict[str, Any],
    print_debug_for_sat: Optional[int] = None
) -> None:
    """
    Updates the visibility table for each satellite against each fixed point.

    Args:
        data_struct: The main simulation data dictionary.
        print_debug_for_sat: If an integer is provided, the `exclusion` function's
                             debug printout will be enabled for that satellite index.
    """
    num_satellites = data_struct['counts']['satellites']
    fixed_points = data_struct['fixedpoints']['position']
    num_fixed_points = len(fixed_points)
    visibility_table = data_struct['fixedpoints']['visibility']

    # Ensure the visibility table has the correct shape
    if visibility_table.shape != (num_fixed_points, num_satellites):
        data_struct['fixedpoints']['visibility'] = np.zeros((num_fixed_points, num_satellites), dtype=int)
        visibility_table = data_struct['fixedpoints']['visibility']

    if num_satellites == 0 or num_fixed_points == 0:
        return

    satellite_positions = data_struct['satellites']['position']

    for i in range(num_satellites):
        sat_pos = satellite_positions[i]
        should_print_debug = (i == print_debug_for_sat)
        for j in range(num_fixed_points):
            fixed_point_pos = fixed_points[j]
            pointing_vector = fixed_point_pos - sat_pos
            data_struct['satellites']['pointing'][i] = pointing_vector
            visibility_table[j, i] = exclusion(data_struct, i, print_debug=should_print_debug)
\end{Verbatim}

\end{document}